\providecommand{\binpie}[2]{\begin{tikzpicture}[baseline=-0.4ex]%
  \fill[gblue!35] (0,0) -- (90:0.5) arc (90:{90-360*#2/#1}:0.5) -- cycle;%
  \foreach \i in {1,...,#1}{\draw[thin] (0,0) -- ({90-360*\i/#1}:0.5);}%
  \draw[thick] (0,0) circle (0.5);\end{tikzpicture}}
\kefalaio{Ρητοί αριθμοί}{Κάθε αριθμός έχει θέση, απόσταση και κατεύθυνση}{%
\begin{tabular}{@{}l@{\hspace{1em}}p{0.82\linewidth}@{}}
3.1 & Η έννοια του ρητού\\[2pt]
3.2 & Απόλυτη τιμή ρητού, σύγκριση και διάταξη ρητών, αντίθετοι ρητοί\\[2pt]
3.3 & Πρόσθεση και αφαίρεση ρητών\\[2pt]
3.4 & Πολλαπλασιασμός ρητών, ιδιότητες του πολλαπλασιασμού, αντίστροφοι ρητοί\\[2pt]
3.5 & Διαίρεση ρητών\\[2pt]
3.6 & Δυνάμεις ρητών με εκθέτη μη μηδενικό φυσικό\\[2pt]
3.7 & Τυποποιημένη μορφή μεγάλων αριθμών\\[2pt]
 & Επανάληψη του κεφαλαίου\\
\end{tabular}}
\paragrafos{Δεκαδικά κλάσματα και δεκαδικοί: ανάμεσα στους ακέραιους}{Α.3.1, «Δεκαδικά κλάσματα – Δεκαδικοί αριθμοί – Διάταξη – Στρογγυλοποίηση»}

\begin{anagke}
Αθλητική μέρα στο καινούργιο σχολείο, και για πρώτη φορά οι δύο γειτονιές τρέχουν στην ίδια πίστα. Η κυρία Αλεξάνδρα κρατάει το χρονόμετρο στα 60 μέτρα και η Πολυξένη γράφει τους χρόνους στο τετράδιο, γιατί στη Θεσσαλονίκη είχε τρέξει σε σύλλογο.

Ο Νίκος τερματίζει πρώτος και η καθηγήτρια φωνάζει: «Εννιά και ογδόντα!». Λίγο μετά τερματίζει ο Βασίλης: «Εννιά και εβδομήντα πέντε». Ο Βασίλης σκύβει το κεφάλι: «Το εβδομήντα πέντε είναι πολύ μεγαλύτερο από το οκτώ. Άρα έκανα πολύ περισσότερο χρόνο από σένα, έχασα με μεγάλη διαφορά». Ο Νίκος διαμαρτύρεται: «Μα έφτασες σχεδόν μαζί μου, με είχες πάνω από την πλάτη μου!». Η Πολυξένη γράφει τους δύο χρόνους, $9{,}8$ και $9{,}75$, και παρατηρεί ότι δεν είναι ίδιου μήκους: «Το ένα έχει ένα ψηφίο μετά την υποδιαστολή και το άλλο δύο. Πώς τα συγκρίνεις;»

Ο Στέλιος προσθέτει το δικό του πρόβλημα: «Με τι μονάδα μετράμε τέτοιους χρόνους; Το δευτερόλεπτο είναι πολύ χοντρή μονάδα. Αν ο χρόνος μου είναι ανάμεσα στο 9 και στο 10, δεν μπορώ να πω ούτε 9 ούτε 10». Η Μυρτώ το είχε πάθει και στο θρανίο: το πλάτος του ήταν ανάμεσα στα 62 και στα 63 εκατοστά.

Οι φυσικοί δεν φτάνουν. Χρειαζόμαστε αριθμούς \emph{ανάμεσα} στους φυσικούς, που να τους κόβουν σε δέκατα, εκατοστά, χιλιοστά, όσο λεπτά χρειάζεται η μέτρηση. Τέτοιους αριθμούς τους ξέρουμε ήδη από τα κλάσματα, και τώρα θα τους γράψουμε με τρόπο που συγκρίνεται εύκολα.
\end{anagke}

\begin{lexi}{δεκαδικός, υποδιαστολή}{δεκαδικός, υποδιαστολή}
Η λέξη \lexiword{δεκαδικός} έρχεται από το «δέκα»: ο αριθμός είναι γραμμένος με βάση το 10, και μετά την υποδιαστολή συνεχίζονται οι θέσεις με δέκατα, εκατοστά, χιλιοστά. Η \lexiword{υποδιαστολή} είναι το κόμμα που χωρίζει τον ακέραιο από το δεκαδικό μέρος· «υπό-διαστέλλω» σημαίνει «χωρίζω από κάτω», και παλιότερα ήταν στο ύψος της βάσης των ψηφίων. Το \lexiword{δεκαδικό κλάσμα} είναι το κλάσμα που έχει παρονομαστή δύναμη του 10: $10,\,100,\,1\,000,\,10\,000,\dots$.
\end{lexi}

\begin{metafrasi}
Τρία δέκατα. & $\dfrac{3}{10}=0{,}3$\\
Εκατόν τριάντα πέντε εκατοστά. & $\dfrac{135}{100}=1{,}35$\\
Δύο ακέραια και επτά χιλιοστά. & $2{,}007=2+\dfrac{7}{1\,000}$\\
Ένα κλάσμα $\dfrac{\kappa}{\nu}$ γίνεται δεκαδικός. & $\kappa:\nu$\\
Οι ίδιοι χρόνοι, με τα ίδια ψηφία. & $9{,}8=9{,}80$\\
Ποιος δεκαδικός είναι μεγαλύτερος; & συγκρίνω ακέραιο, δέκατα, εκατοστά, …\\
Προσεγγίζω στα εκατοστά. & $9{,}7468\approx 9{,}75$\\
\end{metafrasi}

\tmima{Η θεωρία}

\tmimasmall{Δεκαδικά κλάσματα και δεκαδικοί αριθμοί}

Τα κλάσματα $\frac{3}{10}$, $\frac{825}{100}$, $\frac{53}{1\,000}$ και $\frac{1004}{10\,000}$ έχουν παρονομαστές τους $10$, $10^2$, $10^3$ και $10^4$. Είναι δεκαδικά κλάσματα.

\begin{orismos}{δεκαδικό κλάσμα και δεκαδικός αριθμός}
Δεκαδικό κλάσμα λέγεται κάθε κλάσμα με παρονομαστή δύναμη του 10. Κάθε δεκαδικό κλάσμα γράφεται ως \textbf{δεκαδικός αριθμός} με τόσα δεκαδικά ψηφία όσα μηδενικά έχει ο παρονομαστής του:
$$\dfrac{3}{10}=0{,}3,\qquad \dfrac{825}{100}=8{,}25,\qquad \dfrac{53}{1\,000}=0{,}053,\qquad \dfrac{1004}{10\,000}=0{,}1004.$$
\end{orismos}

Για να πάω από το δεκαδικό στο κλάσμα, κάνω το αντίστροφο: γράφω στον αριθμητή όλα τα ψηφία χωρίς την υποδιαστολή και στον παρονομαστή το $1$ με τόσα μηδενικά όσα τα δεκαδικά ψηφία. Άρα $3{,}25=\frac{325}{100}$.

Για να μετατρέψω ένα \emph{οποιοδήποτε} κλάσμα $\frac{\kappa}{\nu}$ σε δεκαδικό, κάνω τη διαίρεση $\kappa:\nu$, βάζοντας υποδιαστολή και μηδενικά στον διαιρετέο όταν τελειώσει το ακέραιο μέρος. Αν η διαίρεση τελειώνει, ο δεκαδικός είναι ακριβής, π.χ. $\frac{3}{8}=0{,}375$. Αν δεν τελειώνει ποτέ, όπως στο $\frac{1}{3}=0{,}333\ldots$, τον γράφουμε κατά προσέγγιση (θα το ξαναδούμε στο κεφάλαιο των ρητών).

\tmimasmall{Οι τάξεις των δεκαδικών ψηφίων}

Δεξιά από την υποδιαστολή οι θέσεις μικραίνουν κατά δέκα, όπως αριστερά μεγαλώνουν:

\begin{center}\small
\renewcommand{\arraystretch}{1.4}\setlength{\tabcolsep}{3pt}
\begin{tabular}{@{}cccccccc@{}}
\toprule
εκατοντάδες & δεκάδες & μονάδες & \textbf{,} & δέκατα & εκατοστά & χιλιοστά & δεκάκις χιλιοστά\\
$100$ & $10$ & $1$ & & $\frac{1}{10}$ & $\frac{1}{100}$ & $\frac{1}{1\,000}$ & $\frac{1}{10\,000}$\\
\bottomrule
\end{tabular}
\end{center}

Το επόμενο ψηφίο λέγεται εκατοντάκις χιλιοστό ($\frac{1}{100\,000}$), μετά εκατομμυριοστό, και ούτω καθεξής. Ο $305{,}0408$ έχει δεκαδικό μέρος $0{,}0408$ και διαβάζεται «τριακόσια πέντε ακέραια και τετρακόσια οκτώ δεκάκις χιλιοστά», γιατί το δεκαδικό μέρος $0{,}0408=\frac{408}{10\,000}$. Μπορούμε να τον αναλύσουμε:
$$305{,}0408=3\cdot100+0\cdot10+5\cdot1+\dfrac{0}{10}+\dfrac{4}{100}+\dfrac{0}{1\,000}+\dfrac{8}{10\,000}.$$

Αν προσθέσω ή σβήσω μηδενικά στο τέλος του δεκαδικού μέρους, ο αριθμός δεν αλλάζει: $9{,}8=9{,}80=9{,}800$, γιατί $\frac{8}{10}=\frac{80}{100}=\frac{800}{1\,000}$. Δεν αλλάζει όμως αν βάλω μηδενικό \emph{πριν} από τα ψηφία του δεκαδικού μέρους: $0{,}5\neq 0{,}05$.

\tmimasmall{Μια ιστορική σημείωση: ο Στέβιν και ο Νάπιερ}

Οι Βαβυλώνιοι και οι αρχαίοι Έλληνες αστρονόμοι χρησιμοποιούσαν κλάσματα με παρονομαστή 60, και οι έμποροι έβρισκαν δύσκολα τους υπολογισμούς με κλάσματα. Ο Φλαμανδός μαθηματικός \textbf{Σιμόν Στέβιν} (1548--1620) κατάλαβε ότι οι υπολογισμοί γίνονται πολύ πιο εύκολα με κλάσματα που έχουν παρονομαστή δύναμη του 10. Το 1585 έγραψε ένα μικρό βιβλίο, τη «Δεκάτη» (\emph{La disme}), για τους τοκογλύφους, τους τοπογράφους και τους εμπόρους, όπου εισήγαγε τον δεκαδικό υπολογισμό. Ο Στέβιν έγραφε τον $34{,}25$ ως $34_{(0)}\,2_{(1)}\,5_{(2)}$: ο δείκτης έδειχνε πόσες θέσεις μετά τη μονάδα βρίσκεται το κάθε ψηφίο. Το κόμμα (ή τελεία) ανάμεσα στον ακέραιο και στο δεκαδικό μέρος το πρότεινε αργότερα ο Σκοτσέζος \textbf{Τζον Νάπιερ} (1550--1617), και η απλότητά του επικράτησε.

\tmimasmall{Σύγκριση δεκαδικών}

\begin{kanonas}{σύγκριση δεκαδικών}
\begin{enumerate}[leftmargin=1.6em,itemsep=0.3ex,topsep=2pt]
\item Συγκρίνω πρώτα τα \textbf{ακέραια μέρη}. Μεγαλύτερο ακέραιο μέρος σημαίνει μεγαλύτερο αριθμό.
\item Αν τα ακέραια μέρη είναι ίσα, συγκρίνω τα δέκατα, μετά τα εκατοστά, μετά τα χιλιοστά, και σταματάω στο πρώτο ψηφίο που διαφέρει.
\item Για να μην μπερδεύομαι, γράφω στον αριθμό με τα λιγότερα δεκαδικά ψηφία μηδενικά στο τέλος, ώστε να έχουν το ίδιο πλήθος.
\end{enumerate}
\end{kanonas}

Έτσι ο χρόνος του Νίκου $9{,}8=9{,}80$ και του Βασίλη $9{,}75$: τα ακέραια ίσα, τα δέκατα $8>7$, άρα $9{,}80>9{,}75$. Ο Βασίλης συνέκρινε τα δεκαδικά μέρη $75$ και $8$ σαν να ήταν φυσικοί αριθμοί. Στην πραγματικότητα έκανε \emph{μικρότερο} χρόνο και κέρδισε, με διαφορά πέντε εκατοστά του δευτερολέπτου.

\tmimasmall{Στρογγυλοποίηση δεκαδικών}

Όταν δεν χρειαζόμαστε όλα τα ψηφία, τα κόβουμε και κρατάμε μια προσέγγιση. Για να στρογγυλοποιήσω σε μια δεδομένη τάξη (π.χ. στα εκατοστά):
\begin{itemize}[leftmargin=1.2em,itemsep=0.3ex,topsep=2pt]
\item κοιτάζω το \textbf{επόμενο} ψηφίο (στα χιλιοστά, αν στρογγυλοποιώ στα εκατοστά)·
\item αν είναι $0,1,2,3,4$ αφήνω το ψηφίο της τάξης όπως είναι και κόβω τα υπόλοιπα·
\item αν είναι $5,6,7,8,9$ αυξάνω το ψηφίο της τάξης κατά $1$ και κόβω τα υπόλοιπα.
\end{itemize}

\tmimasmall{Οι δεκαδικοί στην ευθεία}

Για να βάλω δεκαδικούς στην ευθεία των αριθμών, χωρίζω κάθε μονάδα σε $10$ ίσα μέρη (δέκατα) και, αν χρειάζεται, κάθε δέκατο σε $10$ ίσα μέρη (εκατοστά), και πάει λέγοντας. Δεν υπάρχει όριο: ανάμεσα σε δύο οποιουσδήποτε διαφορετικούς δεκαδικούς υπάρχει πάντα άλλος δεκαδικός, π.χ. ο μέσος τους, και άρα υπάρχουν άπειροι.

\begin{center}
\begin{tikzpicture}[x=1cm,y=1cm]
\draw[thick,-{Stealth}] (-0.3,0)--(10.7,0);
\foreach \i in {0,...,10}{\draw (\i,-0.12)--(\i,0.12);}
\node[below] at (0,-0.15) {$2$};
\node[below] at (5,-0.15) {$2{,}5$};
\node[below] at (10,-0.15) {$3$};
\foreach \i/\t in {1/2{,}1,3/2{,}3,8/2{,}8}{\node[below,font=\small] at (\i,-0.15) {$\t$};}
\fill[gblue] (7.5,0) circle (2.5pt);
\node[above,gblue] at (7.5,0.2) {$2{,}75$};
\fill[gblue] (3,0) circle (2.5pt);
\node[above,gblue] at (3,0.2) {$2{,}3$};
\end{tikzpicture}
\end{center}

Το $2{,}75$ βρίσκεται στη μέση του διαστήματος από το $2{,}7$ ως το $2{,}8$, που αντιστοιχεί στο σημείο $7{,}5$ των δέκατων.

\begin{prosoxi}
\begin{itemize}[leftmargin=1.2em,itemsep=0.3ex,topsep=2pt]
\item Δεν συγκρίνουμε τα δεκαδικά μέρη \emph{σαν φυσικούς}. Το $0{,}9$ είναι μεγαλύτερο από το $0{,}12$, γιατί $0{,}90>0{,}12$, αν και το $9<12$.
\item Το $0{,}50$ ισούται με το $0{,}5$, γιατί το μηδενικό στο τέλος δεν αλλάζει τίποτα, αλλά το $0{,}05$ είναι δέκα φορές μικρότερο.
\item Στρογγυλοποιούμε \emph{μια φορά}, από τον αρχικό αριθμό. Αν στρογγυλοποιήσουμε πρώτα στα εκατοστά και μετά στα δέκατα, μπορεί να βγει λάθος: $2{,}449\approx2{,}4$, αλλά $2{,}449\to2{,}45\to2{,}5$.
\item Το $\frac{2}{5}$ δεν έχει παρονομαστή δύναμη του 10, δεν είναι δηλαδή δεκαδικό κλάσμα, αλλά είναι ισοδύναμο με το δεκαδικό κλάσμα $\frac{4}{10}=0{,}4$.
\end{itemize}
\end{prosoxi}

\tmima{Λυμένα παραδείγματα}

\paradeigma{Ποιος κέρδισε;}
\tixero Οι χρόνοι στα 60 μ. του Νίκου και του Βασίλη είναι $9{,}8$ s και $9{,}75$ s.\par
\tithelo Ποιος έκανε τον μικρότερο χρόνο.\par
\possindeo Ισοσκελίζω τα δεκαδικά ψηφία: $9{,}8=9{,}80$. Τα ακέραια είναι ίσα (9) και στα δέκατα $8>7$, άρα $9{,}80>9{,}75$.\par
\apantisi Ο Βασίλης έκανε μικρότερο χρόνο ($9{,}75$ s) και κέρδισε με διαφορά πέντε εκατοστά του δευτερολέπτου.\par
\elegxos Σε εκατοστά, $980$ εκατοστά και $975$ εκατοστά: το $975$ είναι μικρότερο. Συμφωνεί.

\paradeigma{Από κλάσμα σε δεκαδικό και ανάποδα}
\tixero Τα κλάσματα $\frac{3}{8}$ και $\frac{7}{20}$ και ο δεκαδικός $0{,}625$.\par
\tithelo Να τα γράψω στην άλλη μορφή.\par
\possindeo Για το $\frac{3}{8}$ κάνω τη διαίρεση $3:8$: $3=3{,}000$ και $3000:8=375$, άρα $0{,}375$. Το $\frac{7}{20}$ γίνεται δεκαδικό κλάσμα αν πολλαπλασιάσω αριθμητή και παρονομαστή με $5$: $\frac{35}{100}=0{,}35$. Και ο $0{,}625$ έχει τρία δεκαδικά ψηφία, άρα $\frac{625}{1\,000}$, που απλοποιείται με το $125$ σε $\frac{5}{8}$.\par
\apantisi $\frac{3}{8}=0{,}375$, $\frac{7}{20}=0{,}35$ και $0{,}625=\frac{5}{8}$.\par
\elegxos $0{,}375\cdot 8=3$ και $\frac{5}{8}:1=0{,}625$ (αφού $5:8=0{,}625$). Σωστά.

\paradeigma{Διαβάζω και αναλύω}
\tixero Ο αριθμός $305{,}0408$.\par
\tithelo Να πω τι ψηφίο βρίσκεται σε κάθε θέση και να τον διαβάσω.\par
\possindeo Πριν από την υποδιαστολή: $3$ εκατοντάδες, $0$ δεκάδες, $5$ μονάδες. Μετά: $0$ δέκατα, $4$ εκατοστά, $0$ χιλιοστά, $8$ δεκάκις χιλιοστά. Το δεκαδικό μέρος είναι $0{,}0408=\frac{408}{10\,000}$.\par
\apantisi «Τριακόσια πέντε ακέραια και τετρακόσια οκτώ δεκάκις χιλιοστά».\par
\elegxos $\frac{4}{100}+\frac{8}{10\,000}=\frac{400}{10\,000}+\frac{8}{10\,000}=\frac{408}{10\,000}$. Σωστά.

\paradeigma{Στρογγυλοποιώ σε τάξεις}
\tixero Ο αριθμός $16{,}4963$.\par
\tithelo Να τον στρογγυλοποιήσω στην μονάδα, στα δέκατα, στα εκατοστά και στα χιλιοστά.\par
\possindeo Κοιτάζω κάθε φορά το επόμενο ψηφίο. Μονάδα: επόμενο το $4$ (<5) άρα $16$. Δέκατα: επόμενο το $9$ (≥5) άρα $4+1=5$, δηλαδή $16{,}5$. Εκατοστά: επόμενο το $6$ άρα $9+1=10$, που δίνει $16{,}50$. Χιλιοστά: επόμενο το $3$ άρα $16{,}496$.\par
\apantisi $16$, $16{,}5$, $16{,}50$, $16{,}496$.\par
\elegxos Το σφάλμα στα εκατοστά είναι $16{,}50-16{,}4963=0{,}0037<0{,}005$, δηλαδή μικρότερο από μισό εκατοστό. Σωστό. Το μηδενικό στο $16{,}50$ μένει για να δείχνει ότι μετρήσαμε μέχρι τα εκατοστά.

\paradeigma{Ένας δεκαδικός ανάμεσα}
\tixero Οι αριθμοί $3{,}14$ και $3{,}15$.\par
\tithelo Να βρω έναν δεκαδικό ανάμεσά τους.\par
\possindeo Γράφω τα δύο με τρία δεκαδικά: $3{,}140$ και $3{,}150$. Ανάμεσα είναι όλοι οι $3{,}141,\,3{,}142,\dots,3{,}149$.\par
\apantisi Για παράδειγμα ο $3{,}145$, που είναι ο μέσος τους.\par
\elegxos $3{,}14<3{,}145<3{,}15$. Και αφού μπορώ να ξαναχωρίσω το διάστημα, υπάρχουν άπειροι.

\tmima{Ασκήσεις}

\zoni{Το κατάλαβα}{δεκαδικά κλάσματα, ανάγνωση, σύγκριση}

\ex Γράψε ως δεκαδικούς: α) $\frac{7}{10}$ \quad β) $\frac{43}{100}$ \quad γ) $\frac{9}{1\,000}$ \quad δ) $\frac{2051}{1\,000}$ \quad ε) $\frac{38}{10\,000}$ \quad στ) $\frac{5}{100\,000}$.

\ex Γράψε ως δεκαδικά κλάσματα (με έναν αριθμητή και έναν παρονομαστή): α) $0{,}6$ \quad β) $0{,}07$ \quad γ) $3{,}25$ \quad δ) $0{,}0019$ \quad ε) $12{,}5$ \quad στ) $0{,}400$.

\ex Μετάτρεψε σε δεκαδικούς με διαίρεση ή με ισοδύναμο κλάσμα: α) $\frac{3}{4}$ \quad β) $\frac{7}{8}$ \quad γ) $\frac{9}{20}$ \quad δ) $\frac{11}{25}$ \quad ε) $\frac{5}{16}$ \quad στ) $\frac{3}{5}$.

\ex Γράψε με ψηφία: α) τρία ακέραια και είκοσι πέντε χιλιοστά \quad β) τετρακόσια εβδομήντα δύο εκατοστά \quad γ) δύο δέκατα \quad δ) δεκαπέντε δεκάκις χιλιοστά \quad ε) έξι ακέραια, πέντε εκατοστά και ένα εκατοντάκις χιλιοστό.

\ex Ποιο ψηφίο βρίσκεται στα εκατοστά και ποιο στα δεκάκις χιλιοστά του $24{,}0713$; Πώς διαβάζεται ο αριθμός;

\ex Ανάλυσε τους αριθμούς όπως στη θεωρία: α) $4{,}073$ \quad β) $12{,}408$ \quad γ) $0{,}0905$.

\ex Σωστό ή λάθος; α) $0{,}5=0{,}50$ \quad β) $0{,}5=0{,}05$ \quad γ) $3{,}40=3{,}4$ \quad δ) $0{,}0070=0{,}007$ \quad ε) Ο $12{,}5$ διαβάζεται «δώδεκα και πέντε» \quad στ) $\frac{1}{4}=0{,}4$ \quad ζ) $\frac{25}{100}=0{,}25$.

\ex Σύγκρινε τα ζευγάρια με τα σύμβολα $<,\,=,\,>$: α) $3{,}07$ και $3{,}7$ \quad β) $0{,}58$ και $0{,}6$ \quad γ) $12{,}09$ και $12{,}090$ \quad δ) $5{,}999$ και $6{,}001$ \quad ε) $0{,}3$ και $0{,}29$ \quad στ) $7{,}45$ και $7{,}5$.

\ex Στρογγυλοποίησε τους $7{,}486$ και $0{,}0951$: α) στη μονάδα \quad β) στα δέκατα \quad γ) στα εκατοστά.

\zoni{Το χρησιμοποιώ}{ο στίβος, η ζυγαριά, η πινακίδα}

\ex \textbf{Ο στίβος.} Στα 60 μ. οι χρόνοι ήταν: Νίκος $9{,}38$ s, Βασίλης $9{,}4$ s, Στέλιος $9{,}29$ s, Πολυξένη $9{,}35$ s, Αγγελική $9{,}3$ s, Ηλίας $9{,}41$ s. Γράψε τους μαθητές με τη σειρά που τερμάτισαν, από τον πρώτο.

\ex \textbf{Το άλμα.} Στο άλμα εις μήκος η Ζωή πήδηξε $3{,}08$ μ., η Μυρτώ $2{,}98$ μ., η Φανή $3{,}1$ μ. και η Πολυξένη $3{,}085$ μ. Γράψε τα μήκη από το μεγαλύτερο στο μικρότερο και πες ποια παιδιά πήραν το χρυσό και το ασημένιο μετάλλιο.

\ex \textbf{Η ζυγαριά.} Η Αγγελική ζυγίζει τα μήλα του παππού και βλέπει $2{,}35$ κιλά. Πόσα είναι στρογγυλοποιημένα στο κιλό και πόσα στα δέκατα του κιλού;

\ex \textbf{Η πινακίδα.} Στο βενζινάδικο η πινακίδα γράφει $1{,}879$ € το λίτρο. Ο θείος Λευτέρης λέει ότι «κοντά στα $1{,}9$», και η Πολυξένη ότι «κοντά στα $1{,}88$». Ποια στρογγυλοποίηση είναι ποια, και ποια είναι πιο κοντινή στην πραγματική τιμή;

\ex \textbf{Ανάμεσα.} Βρες τρεις διαφορετικούς δεκαδικούς ανάμεσα στους $2{,}4$ και $2{,}5$. Πόσοι ακόμη υπάρχουν;

\ex \textbf{Στην ευθεία.} Σχεδίασε ένα ευθύγραμμο τμήμα με μήκος $10$ εκ., που θα παριστάνει την απόσταση από το $0$ ως το $1$, και σημείωσε τους $0{,}3$, $0{,}75$, $0{,}9$ και $\frac{1}{2}$.

\ex \textbf{Τελειώνει ή όχι;} Κάνε τις διαιρέσεις (μέχρι να τελειώσουν ή να δεις ότι το ίδιο ψηφίο επαναλαμβάνεται): α) $\frac{3}{20}$ \quad β) $\frac{7}{30}$ \quad γ) $\frac{11}{50}$ \quad δ) $\frac{5}{12}$. Ποια τελειώνουν και ποια όχι;

\zoni{Το σκέφτομαι}{όταν το κόμμα ξεγελάει}

\ex Ο Βασίλης ισχυρίζεται ότι $0{,}9<0{,}12$, γιατί το $9$ είναι μικρότερο από το $12$. Εξήγησε το λάθος του και βρες ένα παράδειγμα με λεφτά ($0{,}90$ € και $0{,}12$ €) που το δείχνει.

\ex Πόσα δεκαδικά ψηφία έχει ο $\frac{7}{16}$; Πώς θα το προέβλεπες χωρίς να κάνεις τη διαίρεση, αν σκεφτείς ότι ο παρονομαστής $16$ είναι $2^4$ και ότι το $10^4=2^4\cdot5^4$;

\ex Με τα ψηφία $3$, $0$ και $8$ (το καθένα μία φορά στο δεκαδικό μέρος) και ακέραιο μέρος το $1$, γράψε τον μικρότερο και τον μεγαλύτερο δεκαδικό.

\ex α) Πόσοι δεκαδικοί, με το πολύ δύο δεκαδικά ψηφία, υπάρχουν ανάμεσα στο $2$ και στο $3$ (χωρίς αυτά); β) Πόσοι από αυτούς έχουν ακριβώς δύο δεκαδικά ψηφία (το τελευταίο όχι μηδέν);

\ex Ο Ηλίας στρογγυλοποιεί τον $2{,}449$ στα δέκατα με δύο βήματα: πρώτα στα εκατοστά ($2{,}45$) και μετά στα δέκατα ($2{,}5$). Ο Νίκος τον στρογγυλοποιεί κατευθείαν. Ποιος έχει δίκιο;

\ex Ο Στέβιν έγραφε τον $7{,}206$ ως $7_{(0)}\,2_{(1)}\,0_{(2)}\,6_{(3)}$. Πώς θα έγραφε τον $18{,}053$; Ποιος αριθμός είναι ο $8_{(0)}\,0_{(1)}\,4_{(2)}\,9_{(3)}$;

\zoni{Για το σπίτι}{μετρώ και καταγράφω}

\ex Μέτρησε το πλάτος του θρανίου σου με ένα μολύβι, με ένα χάρακα και με μετροταινία. Γράψε τα αποτελέσματα και εξήγησε ποια μέτρηση είναι η πιο ακριβής και γιατί χρειάζεται δεκαδικός αριθμός.

\ex Στο σούπερ μάρκετ βρες τρία προϊόντα που η τιμή ανά κιλό ή ανά λίτρο έχει τρία δεκαδικά ψηφία. Στρογγυλοποίησέ τις στο λεπτό του ευρώ και βάλε τες σε σειρά από την πιο ακριβή.

\paragrafos{Απόλυτη τιμή και διάταξη: ποιος είναι δεξιότερα;}{Α.3.2, «Απόλυτη τιμή ρητού, σύγκριση και διάταξη ρητών, αντίθετοι ρητοί»}

\begin{anagke}
Στην ώρα των μαθηματικών η κυρία Αλεξάνδρα ρωτάει: «Ποιος είναι ο επόμενος ρητός μετά το $\frac{2}{5}$;». Ο Βασίλης απαντάει με σιγουριά: «Το $\frac{3}{5}$». Η Ρήνη σηκώνει το χέρι: «Στον χάρακα το $\frac{2}{5}$ είναι στο $0{,}4$ και το $\frac{3}{5}$ στο $0{,}6$. Ανάμεσα υπάρχει το $0{,}5$. Και το $0{,}45$. Και το $0{,}41$!». Ο Νίκος ρωτάει: «Αν δεν υπάρχει επόμενος, πώς τους βάζουμε σε σειρά;».

Λίγο αργότερα ο Ηλίας διαβάζει την πρόγνωση: Ιωάννινα $-3^\circ$, Φλώρινα $-10^\circ$. «Το δέκα είναι μεγαλύτερο από το τρία, άρα η Φλώρινα έχει μεγαλύτερη θερμοκρασία», λέει ο Βασίλης. Η Ζωή διαφωνεί: «Δεν μπορεί να είναι πιο ζεστό το δέκα κάτω από το μηδέν!». Η Αγγελική ζωγραφίζει τον άξονα: «Όσο πιο αριστερά, τόσο πιο κρύο. Το δέκα μετράει πόσο μακριά είμαστε από το μηδέν, και το πρόσημο λέει από ποια μεριά».

Η Πολυξένη το διατυπώνει με δύο ερωτήματα: «Πόσο μακριά είναι κάθε αριθμός από το μηδέν; Και ποιος είναι πιο δεξιά;» Η πρώτη ιδέα θα λέγεται απόλυτη τιμή και η δεύτερη διάταξη. Σε αυτή την παράγραφο θα μάθουμε να συγκρίνουμε οποιουσδήποτε δύο ρητούς: κλάσματα, δεκαδικούς, θετικούς και αρνητικούς.
\end{anagke}

\begin{lexi}{απόλυτη τιμή, αντίθετοι, διάταξη}{απόλυτη τιμή, αντίθετοι, διάταξη}
\lexiword{Απόλυτη} σημαίνει «λυμένη από δεσμούς» (από το «απολύω»): η απόλυτη τιμή είναι ο αριθμός ελεύθερος από το πρόσημό του, δηλαδή μόνο η απόσταση από το $0$. Οι \lexiword{αντίθετοι} βρίσκονται σε «αντίθετη» κατεύθυνση: ίδια απόσταση από το $0$, ο ένας δεξιά και ο άλλος αριστερά. Η \lexiword{διάταξη} είναι η «διά-ταξη», το να βάζεις τους αριθμούς «σε τάξη», δηλαδή σε σειρά.
\end{lexi}

\begin{metafrasi}
Η απόσταση του $-2$ από το $0$ είναι $2$. & $|-2|=2$\\
Η απόσταση του $+\frac{1}{2}$ από το $0$ είναι $\frac{1}{2}$. & $\left|+\frac{1}{2}\right|=\frac{1}{2}$\\
Ο αντίθετος του $+2{,}8$ είναι ο $-2{,}8$. & $-(+2{,}8)=-2{,}8$\\
Το $-3{,}7$ είναι αριστερά του $-2{,}5$. & $-3{,}7<-2{,}5$\\
Το $1{,}6$ είναι δεξιά του $0{,}5$. & $1{,}6>0{,}5$\\
\end{metafrasi}

\tmima{Η θεωρία}

\tmimasmall{Απόλυτη τιμή και αντίθετοι}

Ο ορισμός της απόλυτης τιμής που ξέρουμε από τους ακεραίους ισχύει και για τους ρητούς.

\begin{orismos}{απόλυτη τιμή ρητού}
\textbf{Απόλυτη τιμή} ενός ρητού $\alpha$ είναι η \textbf{απόσταση} του σημείου με τετμημένη $\alpha$ από την αρχή $\mathrm{Ο}$ του άξονα. Συμβολίζεται $|\alpha|$.
\end{orismos}

Για παράδειγμα $|-2|=2$ και $\left|+\frac{1}{2}\right|=\frac{1}{2}$. Η απόλυτη τιμή ενός θετικού είναι ο ίδιος ο αριθμός, ενός αρνητικού ο αντίθετός του και $|0|=0$. Δεν είναι ποτέ αρνητική.

\begin{center}
\akgr[1.8]{-2}{2}{-2/-2,0.5/\frac{1}{2}}
\end{center}

\begin{orismos}{αντίθετοι ρητοί}
\textbf{Αντίθετοι} λέγονται δύο ρητοί που είναι \textbf{ετερόσημοι} και έχουν την \textbf{ίδια απόλυτη τιμή}. Ο αντίθετος του $\alpha$ είναι ο $-\alpha$. Ο αντίθετος του $0$ είναι το $0$ και ισχύει $-(-\alpha)=\alpha$.
\end{orismos}

Οι $+2{,}8$ και $-2{,}8$ είναι αντίθετοι. Το ίδιο και οι $-\frac{1}{7}$ και $+\frac{1}{7}$. Στον άξονα, οι αντίθετοι είναι σημεία συμμετρικά ως προς το $0$.

\tmimasmall{Σύγκριση και διάταξη ρητών}

Από δύο ρητούς \textbf{μεγαλύτερος} είναι αυτός που βρίσκεται \textbf{δεξιότερα} στον άξονα.

\begin{center}
\akgr[1.8]{-4}{2}{-3.7/-3{,}7,-2.5/-2{,}5,0.5/0{,}5,1.6/1{,}6}
\end{center}

Στο σχήμα $1{,}6>0{,}5$ και $-2{,}5>-3{,}7$.

\begin{kanonas}{σύγκριση ρητών με απόλυτες τιμές}
\begin{itemize}[leftmargin=1.2em,itemsep=0.3ex,topsep=2pt]
\item Κάθε θετικός είναι μεγαλύτερος από κάθε αρνητικό. Το $0$ είναι μεγαλύτερο από κάθε αρνητικό και μικρότερο από κάθε θετικό.
\item Από δύο \textbf{θετικούς} μεγαλύτερος είναι αυτός που έχει τη \textbf{μεγαλύτερη} απόλυτη τιμή: $+4{,}3>+4{,}2$, γιατί $|{+4{,}3}|>|{+4{,}2}|$.
\item Από δύο \textbf{αρνητικούς} μεγαλύτερος είναι αυτός που έχει τη \textbf{μικρότερη} απόλυτη τιμή: $-0{,}1>-0{,}2$, γιατί $|{-0{,}1}|=0{,}1<0{,}2=|{-0{,}2}|$.
\end{itemize}
\end{kanonas}

\tmimasmall{Σύγκριση κλασμάτων}

Για κλάσματα με \textbf{θετικούς} όρους ισχύουν τα εξής.

\begin{kanonas}{σύγκριση κλασμάτων}
\begin{itemize}[leftmargin=1.2em,itemsep=0.3ex,topsep=2pt]
\item Από δύο \textbf{ομώνυμα} κλάσματα μεγαλύτερο είναι αυτό με τον \textbf{μεγαλύτερο αριθμητή}: $\frac{2}{5}<\frac{3}{5}$.
\item Από δύο κλάσματα με τον \textbf{ίδιο αριθμητή} μεγαλύτερο είναι αυτό με τον \textbf{μικρότερο παρονομαστή}: $\frac{1}{4}<\frac{1}{2}$.
\item Για να συγκρίνουμε \textbf{ετερώνυμα} κλάσματα, τα κάνουμε \textbf{ομώνυμα} και συγκρίνουμε τους αριθμητές τους. Διαφορετικά, εκτελούμε τις διαιρέσεις και συγκρίνουμε τους δεκαδικούς.
\end{itemize}
\end{kanonas}

\begin{center}
\(\dfrac{2}{5}\)\ \binpie{5}{2}\ \ $<$\ \ \(\dfrac{3}{5}\)\ \binpie{5}{3}\qquad\qquad\(\dfrac{1}{4}\)\ \binpie{4}{1}\ \ $<$\ \ \(\dfrac{1}{2}\)\ \binpie{2}{1}
\end{center}

\tmimasmall{Σύγκριση δεκαδικών}

\begin{kanonas}{σύγκριση δεκαδικών}
\begin{itemize}[leftmargin=1.2em,itemsep=0.3ex,topsep=2pt]
\item Αν έχουν διαφορετικό ακέραιο μέρος, μεγαλύτερος είναι αυτός με το μεγαλύτερο ακέραιο μέρος: $8{,}43>5{,}96$.
\item Αν έχουν το ίδιο ακέραιο μέρος, συγκρίνουμε τα δεκαδικά τους ψηφία \textbf{ένα προς ένα από αριστερά προς τα δεξιά} και μεγαλύτερος είναι αυτός που έχει μεγαλύτερο το πρώτο ψηφίο στο οποίο διαφέρουν: $3{,}132>3{,}1294$, γιατί στα χιλιοστά $3>2$.
\end{itemize}
\end{kanonas}

\tmimasmall{Άπειροι ρητοί ανάμεσα σε δύο ρητούς}

Ανάμεσα σε δύο οποιουσδήποτε διαφορετικούς ρητούς υπάρχει ένα \textbf{άπειρο πλήθος} ρητών. Λέμε ότι οι ρητοί είναι «\textbf{πυκνοί}» στην ευθεία: δεν υπάρχει «επόμενος» ενός ρητού, σε αντίθεση με τους ακεραίους, όπου κάθε ακέραιος έχει επόμενο. Για να βρούμε έναν ρητό ανάμεσα σε δύο άλλους, τους γράφουμε ως ομώνυμα (και αν δεν υπάρχει ακέραιος ανάμεσα στους αριθμητές, πολλαπλασιάζουμε τους όρους με έναν αρκετά μεγάλο αριθμό, π.χ. $10$) ή ως δεκαδικούς.

\begin{prosoxi}
\begin{itemize}[leftmargin=1.2em,itemsep=0.3ex,topsep=2pt]
\item Στους αρνητικούς το «μεγάλο νούμερο» είναι μικρός αριθμός: $-10<-3$ παρά το $10>3$.
\item Ο αντίθετος του $-5$ είναι το $+5$ και όχι το $-5$. Και $|-5|=5$, όχι $-5$.
\item Στα δεκαδικά, ο αριθμός με τα περισσότερα ψηφία δεν είναι μεγαλύτερος: $0{,}3>0{,}2999$. Και ο $-1{,}36$ είναι \emph{μικρότερος} από τον $-1{,}3$.
\item Στα κλάσματα, ο μεγαλύτερος παρονομαστής δεν δίνει μεγαλύτερο κλάσμα: $\frac{1}{8}<\frac{1}{3}$. Τα κομμάτια είναι πιο μικρά.
\item Δεν συγκρίνουμε ετερώνυμα κλάσματα κοιτάζοντας μόνο τους αριθμητές ή μόνο τους παρονομαστές.
\end{itemize}
\end{prosoxi}

\tmima{Λυμένα παραδείγματα}

\paradeigma{Ποια πόλη είναι πιο κρύα}
\tixero Ιωάννινα $-3^\circ$, Φλώρινα $-10^\circ$.\par
\tithelo Ποια θερμοκρασία είναι μικρότερη;\par
\possindeo Και οι δύο είναι αρνητικοί, και μεγαλύτερος είναι αυτός με τη μικρότερη απόλυτη τιμή. Ισχύει $|-10|=10>3=|-3|$.\par
\apantisi $-10<-3$: η Φλώρινα είναι πιο κρύα.\par
\elegxos Στον άξονα το $-10$ είναι αριστερότερα από το $-3$.

\paradeigma{Αντίθετοι και απόλυτες τιμές}
\tixero Οι ρητοί $+4{,}5$, $-\frac{2}{3}$ και $-7$.\par
\tithelo Να βρω τις απόλυτες τιμές και τους αντίθετούς τους.\par
\possindeo Η απόλυτη τιμή είναι η απόσταση από το $0$: $4{,}5$, $\frac{2}{3}$, $7$. Ο αντίθετος του $x$ είναι ο $-x$: $-4{,}5$, $+\frac{2}{3}$, $+7$.\par
\apantisi $|{+4{,}5}|=4{,}5$, $\left|-\frac{2}{3}\right|=\frac{2}{3}$, $|-7|=7$ και αντίθετοι $-4{,}5$, $+\frac{2}{3}$, $+7$.\par
\elegxos Κάθε αριθμός και ο αντίθετός του έχουν την ίδια απόλυτη τιμή.

\paradeigma{Ετερώνυμα κλάσματα}
\tixero Τα κλάσματα $\frac{3}{5}$, $\frac{5}{8}$ και $\frac{7}{12}$.\par
\tithelo Να τα βάλω σε αύξουσα σειρά.\par
\possindeo Τα κάνω ομώνυμα. $\EKP(5,8,12)=120$. $\frac{3}{5}=\frac{72}{120}$, $\frac{5}{8}=\frac{75}{120}$, $\frac{7}{12}=\frac{70}{120}$. Συγκρίνω τους αριθμητές: $70<72<75$.\par
\apantisi $\dfrac{7}{12}<\dfrac{3}{5}<\dfrac{5}{8}$.\par
\elegxos Με δεκαδικούς: $0{,}583\ldots<0{,}6<0{,}625$.

\paradeigma{Αρνητικά κλάσματα και δεκαδικοί}
\tixero Οι αριθμοί $-\frac{5}{2}$, $-0{,}9$, $-4{,}5$, $0{,}1$, $\frac{1}{2}$, $3$.\par
\tithelo Να τους τοποθετήσω στην αριθμογραμμή και να τους συγκρίνω.\par
\possindeo Γράφω τα κλάσματα ως δεκαδικούς: $-2{,}5$ και $0{,}5$. Οι αρνητικοί είναι αριστερά του $0$ και ο μικρότερος είναι αυτός με τη μεγαλύτερη απόλυτη τιμή: $-4{,}5<-2{,}5<-0{,}9$. Μετά οι θετικοί: $0{,}1<0{,}5<3$.\par
\apantisi $-4{,}5<-\dfrac{5}{2}<-0{,}9<0{,}1<\dfrac{1}{2}<3$.\par
\elegxos Η σειρά συμφωνεί με τη σειρά των σημείων στον άξονα, από αριστερά προς τα δεξιά.

\begin{center}
\akgr[1.1]{-5}{4}{-4.5/-4{,}5,-2.5/-\frac{5}{2},-0.9/-0{,}9,0.1/0{,}1,0.5/\frac{1}{2},3/3}
\end{center}

\paradeigma{Ένας ρητός ανάμεσα}
\tixero Οι αριθμοί $\frac{1}{3}$ και $\frac{1}{2}$.\par
\tithelo Να βρω έναν ρητό ανάμεσά τους.\par
\possindeo \emph{1ος τρόπος.} Κάνω τα κλάσματα ομώνυμα: $\frac{1}{3}=\frac{2}{6}$ και $\frac{1}{2}=\frac{3}{6}$. Δεν υπάρχει ακέραιος ανάμεσα στο $2$ και στο $3$, γι' αυτό πολλαπλασιάζω τους όρους επί $10$: $\frac{20}{60}$ και $\frac{30}{60}$. Ανάμεσα βρίσκεται το $\frac{25}{60}$. \emph{2ος τρόπος.} Με δεκαδικούς: $\frac{1}{3}=0{,}333\ldots$ και $\frac{1}{2}=0{,}5$, άρα ανάμεσα βρίσκεται ο $0{,}4$.\par
\apantisi $\dfrac{1}{3}<\dfrac{25}{60}<\dfrac{1}{2}$ και $\dfrac{1}{3}<0{,}4<\dfrac{1}{2}$.\par
\elegxos $\frac{25}{60}=\frac{5}{12}=0{,}41\overline{6}$. Βρίσκεται ανάμεσα στο $0{,}333\ldots$ και στο $0{,}5$.

\paradeigma{Σύγκριση δεκαδικών}
\tixero Οι αριθμοί $9{,}8$ και $9{,}75$, και οι $-1{,}36$ και $-1{,}3$.\par
\tithelo Ποιος είναι μεγαλύτερος σε κάθε ζεύγος;\par
\possindeo Ίδιο ακέραιο μέρος στο πρώτο ζεύγος. Γράφω $9{,}8=9{,}80$ και συγκρίνω τα δέκατα: $8>7$, άρα $9{,}80>9{,}75$. Στο δεύτερο οι αριθμοί είναι αρνητικοί, και $|-1{,}36|=1{,}36>1{,}30=|-1{,}3|$, άρα $-1{,}36<-1{,}3$.\par
\apantisi $9{,}8>9{,}75$ και $-1{,}3>-1{,}36$.\par
\elegxos Αν $9{,}8$ και $9{,}75$ είναι χρόνοι τρεξίματος, ο $9{,}75$ είναι ο καλύτερος.

\tmima{Ασκήσεις}

\zoni{Το κατάλαβα}{απόλυτη τιμή, αντίθετοι, σύγκριση}

\ex Συμπλήρωσε τα κενά. α) Ο αντίθετος του $\alpha$ είναι ο \ldots \quad β) Από δύο θετικούς ρητούς μεγαλύτερος είναι αυτός με τη \ldots απόλυτη τιμή. \quad γ) Από δύο αρνητικούς ρητούς μεγαλύτερος είναι αυτός με τη \ldots απόλυτη τιμή. \quad δ) Από δύο ομώνυμα κλάσματα μεγαλύτερο είναι αυτό με τον μεγαλύτερο \ldots \quad ε) Από δύο κλάσματα με ίδιο αριθμητή μεγαλύτερο είναι αυτό με τον \ldots παρονομαστή.

\ex Συμπλήρωσε τον πίνακα.

\begin{center}
\renewcommand{\arraystretch}{1.6}
\begin{tabular}{|l|c|c|c|c|c|c|}
\hline
Αριθμός & $-3$ & $0{,}35$ & $\frac{1}{9}$ & & & $-\frac{2}{5}$\\
\hline
Αντίθετος & & & & $-0{,}1$ & $\frac{7}{4}$ & \\
\hline
Απόλυτη τιμή & & & & & & \\
\hline
\end{tabular}
\end{center}

\ex Σύγκρινε με ένα από τα $<,\ >,\ =$: α) $\frac{21}{6}\ \ldots\ \frac{12}{6}$ \quad β) $\frac{4}{9}\ \ldots\ \frac{4}{6}$ \quad γ) $\frac{2}{3}\ \ldots\ \frac{7}{8}$ \quad δ) $\frac{1}{100}\ \ldots\ -\frac{2}{5}$ \quad ε) $-\frac{8}{10}\ \ldots\ -1$ \quad στ) $-\frac{8}{3}\ \ldots\ -3$ \quad ζ) $\frac{9}{8}\ \ldots\ 1{,}2$.

\ex Σύγκρινε τους δεκαδικούς: α) $3{,}16\ \ldots\ 3{,}2$ \quad β) $0{,}092\ \ldots\ 0{,}087$ \quad γ) $-1{,}36\ \ldots\ -0{,}01$ \quad δ) $-2{,}69\ \ldots\ 1{,}4$ \quad ε) $-1{,}3\ \ldots\ -1{,}33$ \quad στ) $-1{,}10\ \ldots\ -1{,}1$.

\ex Βάλε σε αύξουσα σειρά τους δεκαδικούς: $1,\ -3,\ 0{,}2,\ 1{,}37,\ -0{,}03,\ 0$.

\ex Βάλε σε αύξουσα σειρά τα κλάσματα: α) $\frac{5}{4},\ \frac{1}{4},\ -\frac{3}{4},\ \frac{3}{4},\ -\frac{1}{4},\ -\frac{5}{4}$ \quad β) $\frac{5}{3},\ \frac{5}{4},\ -\frac{5}{2},\ \frac{5}{8},\ -\frac{5}{7},\ -\frac{5}{3}$.

\ex Δίνονται τα κλάσματα $\frac{9}{4},\ \frac{3}{16},\ \frac{9}{8},\ \frac{15}{16},\ \frac{5}{4},\ \frac{7}{16}$. α) Βρες ποια είναι μικρότερα του $1$ και γράψε τα σε αύξουσα σειρά. β) Βρες ποια είναι μεγαλύτερα του $1$ και γράψε τα σε αύξουσα σειρά. γ) Γράψε και τα έξι σε αύξουσα σειρά.

\ex Σωστό ή λάθος; α) $|\alpha|\ge0$ για κάθε ρητό $\alpha$. \quad β) $|-5|=-5$. \quad γ) Ο αντίθετος του $-4$ είναι ο $4$. \quad δ) $-9>-2$. \quad ε) $\frac{3}{7}<\frac{3}{5}$. \quad στ) Ο αντίθετος του $0$ είναι το $0$.

\zoni{Το χρησιμοποιώ}{θερμόμετρα, ώρες, αθλήματα}

\ex \textbf{Οι θερμοκρασίες.} Στον πίνακα του καιρού: Μέτσοβο $-5{,}5^\circ$, Ιωάννινα $-3^\circ$, Φλώρινα $-10^\circ$, Αθήνα $7^\circ$, Θεσσαλονίκη $-1{,}5^\circ$. Γράψε τις πόλεις από την πιο κρύα στην πιο ζεστή.

\ex \textbf{Οι δραστηριότητες.} Στο κυκλικό διάγραμμα της Αγγελικής ο Ύπνος παίρνει το $\frac{1}{3}$ του εικοσιτετραώρου, το Σχολείο το $\frac{1}{4}$, το Διάβασμα το $\frac{1}{8}$ και η Διασκέδαση το $\frac{1}{6}$. α) Πόσες ώρες διαρκεί η κάθε δραστηριότητα; β) Γράψε τις δραστηριότητες από αυτή που παίρνει τον περισσότερο χρόνο ως αυτή που παίρνει τον λιγότερο. γ) Τι μέρος του εικοσιτετραώρου μένει για όλα τα άλλα;

\ex \textbf{Τα εκατοστά του δευτερολέπτου.} Ο Νίκος τρέχει τα $60$ μέτρα σε $9{,}8$ δευτερόλεπτα και ο Βασίλης σε $9{,}75$. Ο Βασίλης λέει «το $75$ είναι μεγαλύτερο από το $8$, άρα έκανα περισσότερο χρόνο». Ποιος κέρδισε και ποιο είναι το λάθος του Βασίλη;

\ex \textbf{Οι αποστάσεις.} Στον άξονα της οδού το σχολείο είναι στο $0$ (μονάδα $100$ μ.), ο φούρνος στο $+3$ και το βιβλιοπωλείο στο $-2$. Ποιο μαγαζί είναι πιο κοντά στο σχολείο; Πόσο απέχουν μεταξύ τους τα δύο μαγαζιά;

\ex \textbf{Τα υψόμετρα.} Η Νεκρά Θάλασσα είναι στα $-430$ μ., ο Όλυμπος στα $+2\,917$ μ., η επιφάνεια της θάλασσας στο $0$ και ο βυθός της τάφρου των Μαριανών στα $-11\,034$ μ. Γράψε τα σε αύξουσα σειρά υψομέτρου.

\zoni{Το σκέφτομαι}{γιατί δεν υπάρχει επόμενος;}

\ex Ο Βασίλης λέει ότι ο επόμενος ρητός του $\frac{2}{5}$ είναι ο $\frac{3}{5}$. Βρες τρεις ρητούς ανάμεσά τους.

\ex Βρες δύο ρητούς ανάμεσα στους: α) $-0{,}7$ και $-0{,}8$ \quad β) $\frac{3}{10}$ και $\frac{2}{5}$ \quad γ) $\frac{7}{9}$ και $\frac{5}{6}$.

\ex Για ποιους ρητούς ισχύει $|\alpha|=\alpha$ και για ποιους $|\alpha|=-\alpha$; Γιατί ισχύει πάντα $|\alpha|=|-\alpha|$;

\ex Βρες τους δύο διαδοχικούς ακεραίους ανάμεσα στους οποίους βρίσκεται κάθε αριθμός: α) $0{,}3$ \quad β) $-1{,}25$ \quad γ) $\frac{7}{2}$ \quad δ) $-\frac{11}{3}$ \quad ε) $\frac{2024}{2025}$.

\zoni{Για το σπίτι}{ψάξε και σύγκρινε}

\ex Βρες τις θερμοκρασίες πέντε πόλεων του κόσμου και βάλε τις σε σειρά από την πιο κρύα στην πιο ζεστή. Σχεδίασε και τον άξονα.

\ex Ζήτησε από τρία μέλη της οικογένειάς σου να σου πουν τον χρόνο που κοιμούνται (σε ώρες) και γράψε το κλάσμα του εικοσιτετραώρου. Ποιος κοιμάται περισσότερο;

\providecommand{\key}[1]{\mbox{\fbox{\rule{0pt}{1.9ex}\hspace{0.15em}#1\hspace{0.15em}}}\hspace{0.1em}}
\providecommand{\kpl}{\key{$+$}}\providecommand{\kmi}{\key{$-$}}\providecommand{\kpo}{\key{$\times$}}\providecommand{\kdi}{\key{$\div$}}\providecommand{\kis}{\key{$=$}}

\paragrafos{Υπολογιστής τσέπης: ένα εργαλείο, όχι ένα μυαλό}{Α.3.3, «Υπολογισμοί με τη βοήθεια υπολογιστή τσέπης»}

\begin{anagke}
Μετά το παζάρι, το ταμείο έχει γεμίσει αποδείξεις. Ο Νίκος, που κρατάει το ταμείο, αρχίζει να τις προσθέτει με το χέρι και βγάζει $68{,}79$ €. Η Πολυξένη τις ξαναπροσθέτει με το κομπιουτεράκι του κινητού της και βγάζει $77{,}79$ €. «Ένας από τους δύο έχει λάθος», λέει ο Ηλίας. «Το κομπιουτεράκι δεν κάνει λάθος», απαντά ο Βασίλης. «Άρα ο Νίκος».

Η Πολυξένη κοιτάζει την οθόνη της και βλέπει ότι πάτησε κατά λάθος το $23{,}45$ ως $32{,}45$: τα ψηφία αντιστράφηκαν, και η διαφορά στο σύνολο είναι ακριβώς $9$ ευρώ. «Το κομπιουτεράκι έκανε σωστή πράξη με λάθος δεδομένα», παρατηρεί η κυρία Αλεξάνδρα. «Δεν σκέφτεται: κάνει ό,τι του πεις». Και ο Στέλιος, που θέλει να δοκιμάσει το εργαλείο, πατάει $1:3$ και βλέπει $0{,}33333333$ με οκτώ ψηφία. «Χάλασε; Το $\frac{1}{3}$ δεν είναι ακριβώς αυτό!».

Στο κεφάλαιο αυτό η παρέα μαθαίνει να χρησιμοποιεί το κομπιουτεράκι σωστά: πώς πατάμε τις πράξεις, πώς φυλάμε ενδιάμεσα αποτελέσματα στη μνήμη, πώς ξέρουμε αν η οθόνη λέει αλήθεια.
\end{anagke}

\begin{lexi}{υπολογιστής, μνήμη, οθόνη}{υπολογιστής, μνήμη}
Ο \lexiword{υπολογιστής} είναι αυτός που «υπολογίζει», από το «λόγος, λογαριασμός»: το να υπολογίζεις σημαίνει να κάνεις λογαριασμό. Παλιά οι υπολογιστές ήταν άνθρωποι με άβακα. Ο \lexiword{άβακας} ήταν ένας πίνακας με χάντρες. Σήμερα το μηχάνημα έχει δική του \lexiword{μνήμη}, ένα «συρτάρι» όπου φυλάμε έναν αριθμό για να τον ξαναβρούμε αργότερα.
\end{lexi}

\begin{metafrasi}
Πρόσθεσε $12{,}35$ και $7{,}8$. & \key{12,35} \kpl\ \key{7,8} \kis\\
Φύλαξε τον αριθμό της οθόνης στη μνήμη. & \key{M+}\\
Αφαίρεσε από τη μνήμη τον αριθμό της οθόνης. & \key{M$-$}\\
Δείξε μου τι υπάρχει στη μνήμη. & \key{MR}\\
Άδειασε τη μνήμη. & \key{MC}\\
Πρώτα η παρένθεση, μετά το υπόλοιπο. & \key{(}\ldots\key{)} ή ενδιάμεσα αποτελέσματα\\
\end{metafrasi}

\tmima{Η θεωρία}

\tmimasmall{Οι τέσσερις πράξεις}

Όλα τα κομπιουτεράκια εκτελούν τις τέσσερις πράξεις. Για να κάνουμε μια πράξη μεταξύ δύο αριθμών, πληκτρολογούμε τον πρώτο, το σύμβολο της πράξης, τον δεύτερο και μετά το \key{$=$}. Τα σύμβολα είναι \kpl, \kmi, \kpo{} και \kdi. Το κόμμα των δεκαδικών είναι συνήθως τελεία στην οθόνη και στο πλήκτρο. Για παράδειγμα, για το $7{,}28:5{,}2$ πατάμε \key{7.28} \kdi\ \key{5.2} \kis{} και βλέπουμε $1{,}4$.

\tmimasmall{Η μνήμη}

Αν το κομπιουτεράκι έχει μνήμη, έχει τέσσερα πλήκτρα:
\begin{itemize}[leftmargin=1.2em,itemsep=0.3ex,topsep=2pt]
\item \key{MR} (ή \key{MRC}): εμφανίζει στην οθόνη τον αριθμό που βρίσκεται στη μνήμη.
\item \key{MC}: σβήνει το περιεχόμενο της μνήμης.
\item \key{M+}: προσθέτει στη μνήμη τον αριθμό της οθόνης.
\item \key{M$-$}: αφαιρεί από τη μνήμη τον αριθμό της οθόνης.
\end{itemize}
Η μνήμη είναι ό,τι χρειάζεται η παρένθεση του χαρτιού: φυλάμε ένα αποτέλεσμα, κάνουμε άλλο υπολογισμό και μετά τα συνδυάζουμε.

\tmimasmall{Παραστάσεις με πολλές πράξεις}

Δεν εμπιστευόμαστε τυφλά τη σειρά των πλήκτρων. Μερικά απλά κομπιουτεράκια εκτελούν τις πράξεις με τη σειρά που τις πληκτρολογούμε, χωρίς να υπολογίζουν πρώτα τους πολλαπλασιασμούς (τα επιστημονικά την σέβονται). Για παραστάσεις με παρενθέσεις ή μικτές πράξεις, είναι ασφαλές να κάνουμε:
\begin{enumerate}[leftmargin=1.6em,itemsep=0.3ex,topsep=2pt]
\item Μια πράξη τη φορά, με την προτεραιότητα του χαρτιού (παρενθέσεις, δυνάμεις, πολλαπλασιασμός/διαίρεση, μετά πρόσθεση/αφαίρεση).
\item Φύλαξη των αποτελεσμάτων που θα χρειαστούμε αργότερα με \key{M+} ή \key{M$-$}, ή απλή γραφή τους στο τετράδιο.
\end{enumerate}

\tmimasmall{Η οθόνη δείχνει προσέγγιση}

Η οθόνη έχει συγκεκριμένο πλήθος ψηφίων (συνήθως $8$ ή $10$). Όταν το αποτέλεσμα έχει περισσότερα ψηφία, το κομπιουτεράκι το κόβει ή το στρογγυλοποιεί. Γι' αυτό $1:3$ δείχνει $0{,}33333333$ και $2:3$ δείχνει $0{,}66666667$. Άρα $3\cdot0{,}33333333=0{,}99999999$ και όχι $1$. Δεν σημαίνει ότι το κομπιουτεράκι χάλασε, αλλά ότι δείχνει μια προσέγγιση.

\begin{prosoxi}
\begin{itemize}[leftmargin=1.2em,itemsep=0.3ex,topsep=2pt]
\item Το κομπιουτεράκι κάνει σωστά ό,τι του πεις, ακόμα και αν του πεις λάθος. Πριν από κάθε χρήση κάνω μια \textbf{εκτίμηση} με στρογγυλά νούμερα: τόσο περίπου πρέπει να βγει.
\item Δεν ξεχνάμε να σβήσουμε τη μνήμη (\key{MC}) πριν από νέο υπολογισμό, αλλιώς προσθέτουμε σε παλιό αριθμό.
\item Η τιμή που δείχνει η οθόνη μπορεί να είναι στρογγυλοποιημένη. Όταν χρειαζόμαστε ακρίβεια, δεν στρογγυλοποιούμε ενδιάμεσα αποτελέσματα, μόνο το τελικό.
\end{itemize}
\end{prosoxi}

\tmima{Λυμένα παραδείγματα}

\paradeigma{Το σύνολο του ταμείου με τη μνήμη}
\tixero Οι αποδείξεις είναι $12{,}35$, $7{,}8$, $23{,}45$, $9{,}99$ και $15{,}2$ ευρώ.\par
\tithelo Το σύνολο, και ο μέσος όρος ανά απόδειξη.\par
\possindeo Εκτίμηση: $12+8+23+10+15=68$. Πατάω \key{MC}, και για κάθε απόδειξη πληκτρολογώ τον αριθμό και \key{M+}: \key{12.35}\key{M+}\ \key{7.8}\key{M+}\ \key{23.45}\key{M+}\ \key{9.99}\key{M+}\ \key{15.2}\key{M+}. Μετά \key{MR}. Για τον μέσο όρο διαιρώ με $5$: $68{,}79:5=13{,}758$.\par
\apantisi Σύνολο $68{,}79$ €, κατά μέσο όρο $13{,}76$ € (στρογγυλοποιημένο στο λεπτό).\par
\elegxos Η εκτίμηση ($68$) είναι πολύ κοντά στο $68{,}79$, ενώ το $77{,}79$ της Πολυξένης είναι εκτός (η εκτίμηση είναι $68$).

\paradeigma{Σύνθετη παράσταση με μνήμη}
\tixero Η παράσταση $(2{,}4:0{,}6+1{,}5\cdot4)\cdot2{,}5-18{,}9:(0{,}7\cdot5+3{,}5:0{,}5)$.\par
\tithelo Να την υπολογίσω με κομπιουτεράκι που έχει μνήμη.\par
\possindeo Πρώτα η πρώτη παρένθεση: $2{,}4:0{,}6=4$, $1{,}5\cdot4=6$, $4+6=10$, $10\cdot2{,}5=25$. Πατάω \key{M+} (η μνήμη έχει $25$). Μετά η δεύτερη: $0{,}7\cdot5=3{,}5$, $3{,}5:0{,}5=7$, $3{,}5+7=10{,}5$ και $18{,}9:10{,}5=1{,}8$. Πατάω \key{M$-$} (η μνήμη έχει $25-1{,}8$). Τέλος \key{MR}.\par
\apantisi $23{,}2$.\par
\elegxos Εκτίμηση: $(4+6)\cdot2{,}5=25$ και $18{,}9:10{,}5\approx2$, άρα $\approx23$. Συμφωνεί.

\paradeigma{Ο έλεγχος με εκτίμηση}
\tixero Ο Ηλίας πληκτρολογεί $38{,}5\cdot4{,}2$ και η οθόνη δείχνει $16{,}17$.\par
\tithelo Να δω αν είναι σωστό.\par
\possindeo Εκτίμηση: $40\cdot4=160$. Το $16{,}17$ είναι δέκα φορές μικρότερο. Ο Ηλίας πάτησε $3{,}85$ αντί για $38{,}5$ (η υποδιαστολή σε λάθος θέση). Διορθώνω: $38{,}5\cdot4{,}2=161{,}7$.\par
\apantisi $161{,}7$.\par
\elegxos $161{,}7\approx160$. Και $3{,}85\cdot4{,}2=16{,}17$, το λάθος που είδαμε.

\paradeigma{Η οθόνη δείχνει προσέγγιση}
\tixero Το κομπιουτεράκι έχει οθόνη οκτώ ψηφίων.\par
\tithelo Τι δείχνει για το $2:3$ και πόσο βγαίνει το $3\cdot(2:3)$ αν χρησιμοποιήσω την οθόνη;\par
\possindeo $2:3=0{,}666\ldots$ και η οθόνη κρατάει οκτώ ψηφία, στρογγυλοποιώντας το τελευταίο: $0{,}66666667$. Τριπλασιάζω: $3\cdot0{,}66666667=2{,}00000001$.\par
\apantisi Η οθόνη δείχνει $0{,}66666667$, και το γινόμενο με το $3$ βγαίνει $2{,}00000001$ αντί για $2$.\par
\elegxos Η διαφορά $0{,}00000001$ είναι αμελητέα, αλλά δείχνει ότι δουλεύουμε με προσέγγιση.

\paradeigma{Πολλαπλασιασμοί που προστίθενται}
\tixero Στο παζάρι πουλήθηκαν $12$ κέικ των $2{,}8$ €, $8$ μπουκάλια των $1{,}5$ € και $20$ μπισκότα των $0{,}6$ €.\par
\tithelo Τα συνολικά έσοδα.\par
\possindeo Κάθε γινόμενο το προσθέτω στη μνήμη: \key{12}\kpo\key{2.8}\kis\key{M+}\ \key{8}\kpo\key{1.5}\kis\key{M+}\ \key{20}\kpo\key{0.6}\kis\key{M+}\ \key{MR}.\par
\apantisi $33{,}6+12+12=57{,}6$ €.\par
\elegxos Εκτίμηση: $12\cdot3=36$ και $8\cdot1{,}5=12$ και $20\cdot0{,}6=12$, άρα $\approx60$. Συμφωνεί.

\tmima{Ασκήσεις}

\zoni{Το κατάλαβα}{πλήκτρα, μνήμη, εκτίμηση}

\ex Υπολόγισε με κομπιουτεράκι: α) $318{,}76+47{,}092$ \quad β) $1\,000-87{,}654$ \quad γ) $24{,}6\cdot3{,}05$ \quad δ) $91{,}2:0{,}16$.

\ex Με εκτίμηση βρες ποιο αποτέλεσμα είναι σωστό χωρίς να κάνεις την πράξη: α) $48{,}7\cdot5{,}9$ είναι $28{,}733$ ή $287{,}33$ ή $2873{,}3$; β) $19{,}8:4{,}5$ είναι $0{,}44$ ή $4{,}4$ ή $44$;

\ex Γράψε με ποια πλήκτρα θα υπολογίσεις: α) $(4{,}5+2{,}7)\cdot1{,}25$ \quad β) $36:(0{,}9+2{,}1)$. Πόσο βγαίνουν;

\ex Ποια παράσταση υπολογίζεται και πόσο βγαίνει; α) \key{MC} \key{6.5}\kpo\key{4}\kis\key{M+} \key{3.2}\kpo\key{5}\kis\key{M$-$} \key{MR}. \quad β) \key{1.2}\kpl\key{0.8}\kis\kdi\key{0.5}\kis.

\ex Γράψε τα πλήκτρα που χρειάζονται για τον υπολογισμό της $2{,}5\cdot8+1{,}5\cdot6-4{,}25\cdot2$ με χρήση μνήμης. Πόσο είναι το αποτέλεσμα;

\ex Υπολόγισε με κομπιουτεράκι: α) $(3{,}6+4{,}4):0{,}8-2{,}5$ \quad β) $12{,}6:(4{,}2\cdot0{,}5)$ \quad γ) $0{,}25\cdot(10-3{,}2)+1{,}2$.

\ex Σωστό ή λάθος; α) Το κομπιουτεράκι δεν κάνει ποτέ λάθος. \quad β) Η οθόνη μετά το $1:3$ δείχνει την ακριβή τιμή του $\frac{1}{3}$. \quad γ) Το \key{M+} προσθέτει στη μνήμη τον αριθμό της οθόνης. \quad δ) Όλα τα κομπιουτεράκια εκτελούν πρώτα τους πολλαπλασιασμούς. \quad ε) Όταν έχω κομπιουτεράκι δεν χρειάζεται εκτίμηση.

\zoni{Το χρησιμοποιώ}{ταμείο, καύσιμα, μοιρασιά}

\ex \textbf{Το ταμείο.} Στο παζάρι έσοδα ήταν $57{,}6$ € και τα έξοδα (υλικά, χαρτικά) $18{,}35$ €. Πόσα καθαρά κέρδη έχει το ταμείο της εκδρομής; Έλεγξε με εκτίμηση.

\ex \textbf{Η απόδειξη.} Ο Στέλιος πληκτρολόγησε $6{,}75\cdot1{,}2$ και είδε $81$. Εξήγησε τι πάτησε λάθος και βρες το σωστό αποτέλεσμα.

\ex \textbf{Το ταξίδι.} Το αυτοκίνητο του θείου Λευτέρη καταναλώνει $7{,}8$ λίτρα ανά $100$ χλμ. Πόση βενζίνη θα χρειαστεί για μια διαδρομή $312$ χλμ.; Στρογγυλοποίησε στο δέκατο του λίτρου.

\ex \textbf{Η μοιρασιά.} Το ταμείο των $100$ € μοιράζεται εξίσου σε $7$ παιδιά. Τι δείχνει η οθόνη; Πόσα ευρώ παίρνει το καθένα στρογγυλοποιημένα στο λεπτό και πόσα θα έδινε συνολικά το ταμείο με αυτή τη στρογγυλοποίηση;

\ex \textbf{Το ρεύμα.} Μια λάμπα ισχύος $0{,}06$ kW ανάβει $4{,}5$ ώρες την ημέρα για $30$ ημέρες. Με τιμή $0{,}18$ € ανά kWh, πόσο κοστίζει στον μήνα; (Κόστος = ισχύς × ώρες × τιμή.)

\zoni{Το σκέφτομαι}{ο έλεγχος είναι δική μας δουλειά}

\ex Εξήγησε γιατί στο κομπιουτεράκι των $8$ ψηφίων το $3\cdot(1:3)$ δεν δίνει ακριβώς $1$.

\ex Με ένα απλό κομπιουτεράκι, αν πατήσεις \key{2}\kpl\key{3}\kpo\key{4}\kis{} βλέπεις συχνά $20$ και όχι $14$. Ποια σειρά εκτέλεσης ακολούθησε και πώς θα την πατήσεις για να βγει το σωστό;

\ex Χωρίς να κάνεις την πράξη, εξήγησε γιατί το $487\cdot0{,}98$ δεν μπορεί να είναι $4\,773{,}3$ και βρες τη σωστή τάξη μεγέθους.

\ex Πόσες φορές πρέπει να πατήσεις \kpl\ \key{0.1} \kis{} ξεκινώντας από το $0$, για να φτάσεις το $2{,}5$; Με ποια διαίρεση το βρίσκεις αμέσως;

\ex Κάνε τη διαίρεση $1:7$ στο κομπιουτεράκι σου. Γράψε όσα ψηφία φαίνονται και στρογγυλοποίησέ τον στα εκατοστά. Τι μοτίβο βλέπεις στα ψηφία;

\zoni{Για το σπίτι}{δοκίμασε το δικό σου κομπιουτεράκι}

\ex Πάρε τον λογαριασμό ρεύματος ή νερού του σπιτιού, και με το κομπιουτεράκι έλεγξε ένα από τα γινόμενα ή τα αθροίσματα που αναγράφονται. Βγαίνει το ίδιο;

\ex Βρες τις θερμοκρασίες της τελευταίας εβδομάδας και με το κομπιουτεράκι υπολόγισε τον μέσο όρο τους. Στρογγυλοποίησέ τον στο δέκατο.

\paragrafos{Τυποποιημένη μορφή: οι μεγάλοι αριθμοί χωράνε σε μια γραμμή}{Α.3.4, «Τυποποιημένη μορφή μεγάλων αριθμών»}

\begin{anagke}
Η κυρία Αλεξάνδρα ζητάει από την Α3 μια εργασία για το ηλιακό σύστημα. Η Αγγελική γράφει στο τετράδιό της την απόσταση της Γης από τον Ήλιο: $149\,600\,000$ χιλιόμετρα. Ο Ηλίας γράφει τη μάζα της Γης: $5\,976\,000\,000\,000\,000\,000\,000\,000$ κιλά, και το γράφει τρεις φορές, γιατί δεν είναι σίγουρος για τα μηδενικά. Ο Στέλιος δοκιμάζει με τον Ήλιο και χάνεται στο δεύτερο μέτρημα. «Αυτό είναι αδύνατο να διαβαστεί», παραπονιέται η Μυρτώ.

Η Πολυξένη, που έχει το κομπιουτεράκι της, πατάει $149\,600\,000\cdot10$ και η οθόνη δείχνει $1.496\mathrm{E}9$. «Τι είναι αυτό το E9;». Η κυρία Αλεξάνδρα χαμογελάει: «Το κομπιουτεράκι δεν χωράει όλα τα μηδενικά, και τα γράφει με συντομία: $1{,}496$ επί δέκα στην ένατη. Κάνουν ό,τι κάνουν και οι επιστήμονες». Και η Φανή ρωτάει: «Μπορούμε να γράφουμε κι εμείς έτσι, ώστε να μη χάνουμε τα μηδενικά;».
\end{anagke}

\begin{lexi}{τυποποιημένη μορφή}{τυποποιημένη μορφή}
Η λέξη \lexiword{τύπος} σημαίνει καλούπι, πρότυπο. \lexiword{Τυποποιημένη} είναι η μορφή που έχει το ίδιο καλούπι για όλους τους αριθμούς: ένας αριθμός από $1$ μέχρι $9{,}99\ldots$ επί μια δύναμη του $10$. Οι επιστήμονες την χρησιμοποιούν γιατί δείχνει αμέσως το \emph{μέγεθος} του αριθμού (πόσα μηδενικά) και τα σημαντικά του ψηφία.
\end{lexi}

\begin{metafrasi}
Εκατόν σαράντα εννέα εκατομμύρια εξακόσιες χιλιάδες. & $149\,600\,000=1{,}496\cdot10^8$\\
Τόσα μηδενικά όσο ο εκθέτης. & $10^8=100\,000\,000$ (το $1$ και $8$ μηδενικά)\\
Ο $4{,}56\cdot10^{10}$ γραμμένος κανονικά. & $45\,600\,000\,000$\\
Στην οθόνη του κομπιουτεράκι. & $1.496\mathrm{E}8$ σημαίνει $1{,}496\cdot10^8$\\
Σύγκριση: μεγαλύτερος εκθέτης. & $8{,}2\cdot10^5>9{,}1\cdot10^4$\\
\end{metafrasi}

\tmima{Η θεωρία}

\begin{orismos}{τυποποιημένη μορφή}
Ένας (μεγάλος) αριθμός είναι γραμμένος σε \textbf{τυποποιημένη μορφή} όταν είναι γραμμένος ως $\alpha\cdot10^{\nu}$, όπου $\nu$ φυσικός αριθμός και ο $\alpha$ είναι ένας δεκαδικός (ή φυσικός) με ακέραιο μέρος ένα ψηφίο, δηλαδή $1\leq\alpha<10$.
\end{orismos}

Για παράδειγμα
$$3\,140\,000\,000\,000\,000\,000=3{,}14\cdot10^{18},\qquad 234\,000\,000\,000\,000\,000=2{,}34\cdot10^{17}.$$
Και αντίστροφα, ο $5{,}21\cdot10^5$ είναι ο $521\,000$ και ο $2\cdot10^3$ ο $2\,000$.

\tmimasmall{Πώς γράφω έναν αριθμό σε τυποποιημένη μορφή}

\begin{kanonas}{τυποποίηση}
\begin{enumerate}[leftmargin=1.6em,itemsep=0.3ex,topsep=2pt]
\item Μετακινώ την υποδιαστολή (στους ακεραίους είναι στο τέλος) προς τα αριστερά, μέχρι να μείνει \textbf{ένα μη μηδενικό ψηφίο} πριν από αυτήν. Έτσι βρίσκω τον $\alpha$ (με $1\leq\alpha<10$).
\item Ο εκθέτης $\nu$ είναι το πλήθος των θέσεων που μετακίνησα.
\item Γράφω $\alpha\cdot10^{\nu}$. Αφαιρώ τα μηδενικά που περισσεύουν στο τέλος του $\alpha$.
\end{enumerate}
\end{kanonas}

Για να πάω πίσω, από την τυποποιημένη στη δεκαδική μορφή, πολλαπλασιάζω τον $\alpha$ επί $10^{\nu}$: μετακινώ την υποδιαστολή του $\alpha$ δεξιά κατά $\nu$ θέσεις, βάζοντας μηδενικά όπου χρειάζεται. Έτσι $7{,}4\cdot10^3=7\,400$.

\tmimasmall{Σύγκριση και πράξεις}

Για να συγκρίνω δύο αριθμούς σε τυποποιημένη μορφή, κοιτάζω πρώτα τον \textbf{εκθέτη}: μεγαλύτερος εκθέτης σημαίνει μεγαλύτερο αριθμό, επειδή έχει περισσότερα ψηφία. Αν οι εκθέτες είναι ίσοι, συγκρίνω τους $\alpha$.

Για να πολλαπλασιάσω, πολλαπλασιάζω τους $\alpha$ μεταξύ τους και προσθέτω τους εκθέτες του $10$: $(2{,}5\cdot10^4)\cdot(4\cdot10^3)=10\cdot10^7=1\cdot10^8$. Προσοχή: αν το γινόμενο των $\alpha$ βγει $10$ ή περισσότερο, πρέπει να ξαναγράψω το αποτέλεσμα σε τυποποιημένη μορφή.

\begin{prosoxi}
\begin{itemize}[leftmargin=1.2em,itemsep=0.3ex,topsep=2pt]
\item Ο $\alpha$ είναι πάντα μικρότερος από το $10$ και τουλάχιστον $1$. Οι $52\cdot10^2$, $0{,}52\cdot10^4$ και $10\cdot10^5$ δεν είναι τυποποιημένοι (αν και ισούνται με $5{,}2\cdot10^3$, $5{,}2\cdot10^3$ και $1\cdot10^6$ αντίστοιχα).
\item Ο εκθέτης δεν είναι το πλήθος των ψηφίων αλλά το πλήθος των ψηφίων μείον ένα. Ο $45\,600\,000\,000$ έχει $11$ ψηφία και γράφεται $4{,}56\cdot10^{10}$.
\item Δεν συγκρίνω τους $\alpha$ όταν διαφέρουν οι εκθέτες: $9{,}9\cdot10^3<1\cdot10^4$, αν και το $9{,}9>1$.
\end{itemize}
\end{prosoxi}

\tmima{Λυμένα παραδείγματα}

\paradeigma{Η απόσταση Γη -- Ήλιος}
\tixero Η Γη απέχει από τον Ήλιο $149\,600\,000$ χλμ.\par
\tithelo Να το γράψω σε τυποποιημένη μορφή.\par
\possindeo Μετακινώ την υποδιαστολή (που είναι στο τέλος) προς τα αριστερά ώσπου να μείνει ένα ψηφίο πριν από αυτήν: $1{,}49600000$. Μετακίνησα $8$ θέσεις. Αφαιρώ τα τελευταία μηδενικά.\par
\apantisi $1{,}496\cdot10^8$ χλμ.\par
\elegxos $1{,}496\cdot100\,000\,000=149\,600\,000$. Και ο αριθμός έχει $9$ ψηφία, εκθέτης $9-1=8$.

\paradeigma{Από τυποποιημένη σε δεκαδική}
\tixero Οι αριθμοί $7{,}03\cdot10^6$ και $4{,}56\cdot10^{10}$.\par
\tithelo Να τους γράψω κανονικά.\par
\possindeo Ο $7{,}03\cdot10^6$: μετακινώ την υποδιαστολή $6$ θέσεις δεξιά, $7\,030\,000$. Ο $4{,}56\cdot10^{10}$: $10$ θέσεις δεξιά, $45\,600\,000\,000$.\par
\apantisi $7\,030\,000$ και $45\,600\,000\,000$.\par
\elegxos Στον πρώτο υπάρχουν $7$ ψηφία (εκθέτης $6$) και στον δεύτερο $11$ ψηφία (εκθέτης $10$).

\paradeigma{Ποιος είναι μεγαλύτερος;}
\tixero Οι αριθμοί $8{,}2\cdot10^5$, $9{,}1\cdot10^4$, $3{,}5\cdot10^7$ και $2{,}9\cdot10^7$.\par
\tithelo Να τους βάλω σε αύξουσα σειρά.\par
\possindeo Κοιτάζω πρώτα τους εκθέτες: $4<5<7$. Ο $9{,}1\cdot10^4$ είναι ο μικρότερος και ο $8{,}2\cdot10^5$ ακολουθεί. Οι άλλοι δύο έχουν εκθέτη $7$, και $2{,}9<3{,}5$.\par
\apantisi $9{,}1\cdot10^4<8{,}2\cdot10^5<2{,}9\cdot10^7<3{,}5\cdot10^7$.\par
\elegxos Δεκαδικά: $91\,000<820\,000<29\,000\,000<35\,000\,000$.

\paradeigma{Πόση ώρα ταξιδεύει το φως;}
\tixero Το φως ταξιδεύει με $3\cdot10^5$ χλμ. ανά δευτερόλεπτο και ο Ήλιος απέχει $1{,}496\cdot10^8$ χλμ.\par
\tithelo Σε πόσο χρόνο το φως φτάνει από τον Ήλιο στη Γη.\par
\possindeo Χρόνος = απόσταση : ταχύτητα $=\dfrac{1{,}496\cdot10^8}{3\cdot10^5}=\dfrac{1{,}496}{3}\cdot10^3\approx0{,}4987\cdot10^3\approx499$ δευτερόλεπτα. Στα λεπτά: $499:60\approx8{,}3$.\par
\apantisi Περίπου $8{,}3$ λεπτά.\par
\elegxos Σε $8{,}3$ λεπτά, δηλαδή περίπου $499$ δευτερόλεπτα, το φως διανύει $499\cdot300\,000\approx1{,}5\cdot10^8$ χλμ. Συμφωνεί.

\paradeigma{Γινόμενο σε τυποποιημένη μορφή}
\tixero Οι αριθμοί $2{,}5\cdot10^4$ και $4\cdot10^3$.\par
\tithelo Το γινόμενό τους σε τυποποιημένη μορφή.\par
\possindeo Πολλαπλασιάζω τους $\alpha$: $2{,}5\cdot4=10$. Προσθέτω τους εκθέτες: $4+3=7$. Έχω $10\cdot10^7$. Το $10$ δεν είναι μικρότερο του $10$, γι' αυτό γράφω $10\cdot10^7=10^8$.\par
\apantisi $1\cdot10^8$.\par
\elegxos $25\,000\cdot4\,000=100\,000\,000=10^8$.

\tmima{Ασκήσεις}

\zoni{Το κατάλαβα}{τυποποίηση, αποτυποποίηση, σύγκριση}

\ex Γράψε σε τυποποιημένη μορφή: α) $72\,000$ \quad β) $1\,250\,000$ \quad γ) $903\,000\,000$ \quad δ) $6\,000\,000\,000$ \quad ε) $48\,700$ \quad στ) $15\,000\,000\,000\,000$.

\ex Γράψε με δεκαδική μορφή: α) $2{,}7\cdot10^5$ \quad β) $5{,}06\cdot10^3$ \quad γ) $1{,}9\cdot10^7$ \quad δ) $8\cdot10^4$ \quad ε) $3{,}125\cdot10^6$ \quad στ) $4{,}008\cdot10^9$.

\ex Ποιοι από τους αριθμούς είναι σε τυποποιημένη μορφή; α) $5{,}2\cdot10^3$ \quad β) $52\cdot10^2$ \quad γ) $0{,}52\cdot10^4$ \quad δ) $10\cdot10^5$ \quad ε) $9{,}99\cdot10^1$.

\ex Γράψε σε τυποποιημένη μορφή: α) $24\cdot10^5$ \quad β) $0{,}35\cdot10^7$ \quad γ) $150\cdot10^3$ \quad δ) $0{,}008\cdot10^9$.

\ex Σύγκρινε με τα $<$, $=$, $>$: α) $8{,}2\cdot10^5$ και $9{,}1\cdot10^4$ \quad β) $3{,}5\cdot10^7$ και $2{,}9\cdot10^7$ \quad γ) $6\cdot10^9$ και $5{,}99\cdot10^9$ \quad δ) $4{,}4\cdot10^3$ και $4{,}4\cdot10^5$.

\ex Γράψε σε αύξουσα σειρά: $3{,}2\cdot10^6$, $9\cdot10^5$, $1{,}1\cdot10^7$, $4{,}5\cdot10^6$, $7{,}9\cdot10^5$.

\ex Η Πολυξένη βλέπει στην οθόνη του κομπιουτεράκι: α) $4.5\mathrm{E}7$ \quad β) $2.5\mathrm{E}10$. Ποιοι είναι οι αριθμοί; Πόσο δείχνει η οθόνη αν πολλαπλασιάσει τον $12\,000\,000$ με τον $3\,000\,000$;

\zoni{Το χρησιμοποιώ}{ηλιακό σύστημα και ανθρώπινο σώμα}

\ex \textbf{Σελήνη και Ήλιος.} Η Σελήνη απέχει $384\,400$ χλμ. από τη Γη και ο Ήλιος $149\,600\,000$ χλμ. Γράψε τις αποστάσεις σε τυποποιημένη μορφή και πες περίπου πόσες φορές μακρύτερα είναι ο Ήλιος.

\ex \textbf{Το έτος φωτός.} Ένα έτος φωτός είναι $9\,461\,000\,000\,000$ χλμ. Γράψε το σε τυποποιημένη μορφή. Ο κοντινότερος αστέρας, ο Εγγύτατος του Κενταύρου, απέχει περίπου $4{,}3$ έτη φωτός. Πόσα χιλιόμετρα είναι αυτά, σε τυποποιημένη μορφή, με δύο δεκαδικά;

\ex \textbf{Το φως του Ήλιου.} Η Φανή ρωτάει σε πόσο χρόνο φτάνει στη Γη το φως, αν ταξιδεύει με $3\cdot10^5$ χλμ. το δευτερόλεπτο και ο Ήλιος απέχει $1{,}496\cdot10^8$ χλμ.; Απάντησε σε δευτερόλεπτα και σε λεπτά.

\ex \textbf{Πληθυσμοί.} Ο πληθυσμός της Γης είναι περίπου $8\,000\,000\,000$ άνθρωποι και της Ελλάδας περίπου $10\,400\,000$. Γράψε τους σε τυποποιημένη μορφή και βρες περίπου πόσες φορές μεγαλύτερος είναι ο πρώτος.

\ex \textbf{Οι μάζες.} Η μάζα της Γης είναι $5{,}976\cdot10^{24}$ κιλά και του Ήλιου $1{,}983\cdot10^{30}$ κιλά. Γράψε τη μάζα της Γης σε δεκαδική μορφή και βρες περίπου πόσες φορές είναι ο Ήλιος βαρύτερος.

\ex \textbf{Η μνήμη του κινητού.} Η μνήμη του κινητού της Μυρτώς είναι $128$ GB, δηλαδή $1{,}28\cdot10^{11}$ bytes. Μια φωτογραφία έχει περίπου $4$ MB, δηλαδή $4\cdot10^6$ bytes. Πόσες φωτογραφίες χωράνε;

\zoni{Το σκέφτομαι}{οι κανόνες, όταν σπάνε}

\ex Γιατί ο $10\cdot10^5$ δεν είναι σε τυποποιημένη μορφή; Πώς γράφεται σωστά;

\ex Πόσα ψηφία έχει ο $4{,}56\cdot10^{10}$ όταν γραφεί κανονικά; Και ο $1\cdot10^{\nu}$; Εξήγησε πώς το βρίσκεις από τον εκθέτη.

\ex Υπολόγισε τα $(3\cdot10^4)\cdot(2\cdot10^5)$ και $(4\cdot10^3)\cdot(5\cdot10^6)$. Βάλε τα αποτελέσματα σε τυποποιημένη μορφή και διατύπωσε τον κανόνα του πολλαπλασιασμού.

\ex Ο Βασίλης λέει ότι $9{,}9\cdot10^3>1\cdot10^4$, γιατί $9{,}9>1$. Έχει δίκιο;

\ex Αν μετράμε αριθμούς φωναχτά, ένα νούμερο ανά δευτερόλεπτο, πόσα χρόνια χρειαζόμαστε για να φτάσουμε στο $10^{15}$; (Ένα έτος έχει περίπου $3{,}15\cdot10^7$ δευτερόλεπτα.)

\zoni{Για το σπίτι}{ψάξε τους μεγάλους αριθμούς του κόσμου}

\ex Βρες από αξιόπιστη πηγή τη μάζα του Ηλίου, το πλήθος των αστεριών του γαλαξία μας και την ακτίνα της Γης. Γράψε τους σε δεκαδική και σε τυποποιημένη μορφή.

\ex Φτιάξε πέντε δικούς σου αριθμούς, γράψε τους σε τυποποιημένη μορφή και ζήτησε από έναν συμμαθητή να βρει τους αρχικούς. Έλεγξε τις απαντήσεις του.

\paragrafos{Διαίρεση ρητών: πόσες φορές χωράει;}{Α.3.5, «Διαίρεση ρητών»}

\begin{anagke}
Ο παππούς Γιάννης έχει $450$ λίτρα ελαιόλαδο και θέλει να το βάλει σε δοχεία που χωράνε από $1\frac{1}{2}$ λίτρα. «Πόσα δοχεία θα χρειαστώ;» ρωτάει. Ο Ηλίας γράφει $450:1\frac{1}{2}$ και κολλάει: «Πώς διαιρώ με μεικτό;». Ο Στέλιος προσπαθεί να εφαρμόσει έναν κανόνα που μόλις έμαθε: «Διαιρώ τους αριθμητές μεταξύ τους και τους παρονομαστές μεταξύ τους». Τον σταματά η Ρήνη: «Σε ένα δοχείο χωράνε λίγο παραπάνω από ένα λίτρο. Τα δοχεία δεν μπορεί να είναι περισσότερα από $450$». Η Ζωή ρωτάει «πόσες φορές χωράει;». Μετράνε: κάθε $3$ λίτρα γεμίζουν $2$ δοχεία, άρα τα $450$ λίτρα γεμίζουν $300$ δοχεία.

Η Πολυξένη παρατηρεί ότι το $300$ προκύπτει και ως $450\cdot\frac{2}{3}$. «Το να διαιρέσεις με το $\frac{3}{2}$ είναι το ίδιο με το να πολλαπλασιάσεις με το $\frac{2}{3}$, τον αντίστροφό του;» Η Αγγελική προσθέτει ένα δεύτερο πρόβλημα: «Στο ταμείο της εκδρομής έλειπαν $24$ ευρώ. Αν χωριστούν σε $6$ παιδιά, πόσο έλειπε σε καθένα;» Και ο Νίκος: «Και αν ο διαιρέτης είναι δεκαδικός, π.χ. $3{,}6:0{,}45$;».

Σε αυτή την παράγραφο θα μάθουμε να διαιρούμε κάθε είδους ρητούς: ακεραίους, κλάσματα, δεκαδικούς και θετικούς και αρνητικούς.
\end{anagke}

\begin{lexi}{διαίρεση, πηλίκο, λόγος}{διαίρεση, πηλίκο, λόγος}
Η \lexiword{διαίρεση} είναι το «δι-αιρώ», «χωρίζω σε μέρη». Ο \lexiword{διαιρετέος} είναι αυτός που μοιράζεται, ο \lexiword{διαιρέτης} αυτός που μοιράζει και το \lexiword{πηλίκο} (από το «πηλίκος», «πόσος;») το αποτέλεσμα. Το πηλίκο $\alpha:\beta$ ή $\frac{\alpha}{\beta}$ λέγεται και \lexiword{λόγος} του $\alpha$ προς το $\beta$. Ένα κλάσμα που έχει μέσα του κλάσμα λέγεται \lexiword{σύνθετο}, «φτιαγμένο από πολλά κομμάτια».
\end{lexi}

\begin{metafrasi}
Πόσες φορές χωράει το $2$ στο $6$; & $6:2=3$ γιατί $3\cdot2=6$\\
Ομόσημοι. & πηλίκο θετικό\\
Ετερόσημοι. & πηλίκο αρνητικό\\
Διαιρώ και πολλαπλασιάζω τον αντίστροφο. & $\alpha:\beta=\alpha\cdot\dfrac{1}{\beta}$\\
Κλάσμα που έχει κλάσμα μέσα του. & $\dfrac{\frac{\alpha}{\beta}}{\frac{\gamma}{\delta}}$ (σύνθετο)\\
\end{metafrasi}

\tmima{Η θεωρία}

\begin{orismos}{πηλίκο ρητών}
Το πηλίκο της διαίρεσης $\alpha:\beta$ ή $\dfrac{\alpha}{\beta}$ ($\beta\neq0$) λέγεται και \textbf{λόγος} του $\alpha$ προς το $\beta$. Αν $\alpha:\beta=\gamma$, τότε ο $\gamma$ είναι ο μοναδικός αριθμός για τον οποίο ισχύει $\gamma\cdot\beta=\alpha$.
\end{orismos}

\begin{kanonas}{πρόσημο του πηλίκου}
Για να διαιρέσουμε δύο ρητούς, διαιρούμε τις \textbf{απόλυτες τιμές} τους και στο πηλίκο βάζουμε το πρόσημο $+$, αν είναι \textbf{ομόσημοι}, και το πρόσημο $-$, αν είναι \textbf{ετερόσημοι}.
\end{kanonas}

Για παράδειγμα $(+6):(+2)=+3$, γιατί $(+3)\cdot(+2)=+6$. Επίσης $(-8):(+4)=-2$, γιατί $(-2)\cdot(+4)=-8$, και $(-12):(-3)=+4$, γιατί $(+4)\cdot(-3)=-12$.

\begin{kanonas}{διαίρεση με τον αντίστροφο}
Για να \textbf{διαιρέσουμε} δύο ρητούς αρκεί να \textbf{πολλαπλασιάσουμε} τον διαιρετέο με τον \textbf{αντίστροφο} του διαιρέτη:
$$\alpha:\beta=\alpha\cdot\dfrac{1}{\beta}=\dfrac{\alpha}{\beta}\qquad\text{και}\qquad\dfrac{\alpha}{\beta}:\dfrac{\gamma}{\delta}=\dfrac{\alpha}{\beta}\cdot\dfrac{\delta}{\gamma}.$$
\end{kanonas}

Για παράδειγμα $\frac{2}{3}:\frac{4}{5}=\frac{2}{3}\cdot\frac{5}{4}=\frac{10}{12}=\frac{5}{6}$ και $\left(+\frac{2}{5}\right):\left(-\frac{1}{2}\right)=-\frac{2}{5}\cdot\frac{2}{1}=-\frac{4}{5}$.

\tmimasmall{Σύνθετα κλάσματα}

Ένα κλάσμα, του οποίου ένας τουλάχιστον όρος είναι κλάσμα, λέγεται \textbf{σύνθετο}. Η γραμμή του κλάσματος σημαίνει διαίρεση, άρα για να το κάνουμε \textbf{απλό} εκτελούμε τη διαίρεση:
$$\dfrac{\frac{\alpha}{\beta}}{\frac{\gamma}{\delta}}=\dfrac{\alpha}{\beta}:\dfrac{\gamma}{\delta}=\dfrac{\alpha\cdot\delta}{\beta\cdot\gamma}.$$
Αν ο παρονομαστής ή ο αριθμητής είναι ακέραιος, τον γράφουμε με παρονομαστή $1$. Αν υπάρχουν πρόσημα στους όρους, τα βγάζουμε έξω από τα κλάσματα: $\frac{-1}{2}$, $\frac{3}{-4}$ και $-\frac{3}{4}$ είναι ο ίδιος τύπος αριθμών.

\tmimasmall{Διαίρεση δεκαδικών}

Για να διαιρέσουμε δεκαδικό με δεκαδικό, πολλαπλασιάζουμε \textbf{και τον διαιρετέο και τον διαιρέτη} με την ίδια κατάλληλη δύναμη του $10$ ($10$, $100$, \dots) ώστε ο διαιρέτης να γίνει φυσικός αριθμός. Το πηλίκο δεν αλλάζει. Έτσι $4{,}2:1{,}75=420:175=2{,}4$.

\begin{kanonas}{διαίρεση με $10,\ 100,\ 1\,000,\dots$ και με $0{,}1,\ 0{,}01,\dots$}
Όταν διαιρούμε έναν δεκαδικό με $10,\ 100,\ 1\,000,\dots$ μεταφέρουμε την υποδιαστολή τόσες θέσεις \textbf{προς τα αριστερά} όσα είναι τα μηδενικά. Όταν τον διαιρούμε με $0{,}1,\ 0{,}01,\ 0{,}001,\dots$ μεταφέρουμε την υποδιαστολή τόσες θέσεις \textbf{προς τα δεξιά} όσα είναι τα δεκαδικά ψηφία του διαιρέτη.
\end{kanonas}

Για παράδειγμα $12{,}68:100=0{,}1268$ και $12{,}68:0{,}01=1\,268$.

\begin{prosoxi}
\begin{itemize}[leftmargin=1.2em,itemsep=0.3ex,topsep=2pt]
\item Δεν διαιρούμε ποτέ με το $0$. Αν $5:0=x$ τότε θα έπρεπε $x\cdot0=5$, που δεν ισχύει για κανέναν αριθμό. Το $0:5=0$ είναι σωστό.
\item Η διαίρεση δεν «μικραίνει» πάντα: $3:\frac{1}{2}=6>3$. Διαιρώντας με αριθμό μικρότερο του $1$, το πηλίκο μεγαλώνει.
\item Δεν αντιστρέφουμε και τα δύο κλάσματα: $\frac{2}{3}:\frac{4}{5}=\frac{2}{3}\cdot\frac{5}{4}$ και όχι $\frac{3}{2}\cdot\frac{5}{4}$. Αντιστρέφουμε μόνο τον \emph{διαιρέτη}.
\item Η διαίρεση δεν είναι αντιμεταθετική: $6:2=3$ ενώ $2:6=\frac{1}{3}$.
\item Στη διαίρεση με $0{,}1$ η υποδιαστολή πάει \emph{δεξιά} (το πηλίκο μεγαλώνει) και όχι αριστερά.
\end{itemize}
\end{prosoxi}

\tmima{Λυμένα παραδείγματα}

\paradeigma{Τα δοχεία του παππού Γιάννη}
\tixero $450$ λίτρα λάδι και δοχεία των $1\frac{1}{2}=\frac{3}{2}$ λίτρων.\par
\tithelo Πόσα δοχεία θα γεμίσουν;\par
\possindeo Διαιρώ: $450:\frac{3}{2}=450\cdot\frac{2}{3}=\frac{900}{3}=300$.\par
\apantisi $300$ δοχεία.\par
\elegxos $300\cdot\frac{3}{2}=450$.

\paradeigma{Πρόσημα στη διαίρεση}
\tixero Οι διαιρέσεις $(+6):(+2)$, $(-8):(+4)$, $(-12):(-3)$ και $\left(-\frac{2}{5}\right):\left(-\frac{3}{10}\right)$.\par
\tithelo Να βρω τα πηλίκα.\par
\possindeo Διαιρώ τις απόλυτες τιμές και βάζω πρόσημο: ομόσημοι $+$, ετερόσημοι $-$. Για το τελευταίο: $\frac{2}{5}\cdot\frac{10}{3}=\frac{20}{15}=\frac{4}{3}$.\par
\apantisi $+3$, $-2$, $+4$, $+\dfrac{4}{3}$.\par
\elegxos Για κάθε πηλίκο ο πολλαπλασιασμός με τον διαιρέτη δίνει τον διαιρετέο: $(-2)\cdot(+4)=-8$.

\paradeigma{Διαίρεση κλασμάτων}
\tixero Οι $\frac{2}{3}:\frac{4}{5}$ και $\left(+\frac{2}{5}\right):\left(-\frac{1}{2}\right)$.\par
\tithelo Να τις υπολογίσω.\par
\possindeo Πολλαπλασιάζω με τον αντίστροφο του διαιρέτη. Για το πρώτο ο αντίστροφος του $\frac{4}{5}$ είναι ο $\frac{5}{4}$. Για το δεύτερο ο αντίστροφος του $-\frac{1}{2}$ είναι ο $-2$.\par
\apantisi $\dfrac{2}{3}\cdot\dfrac{5}{4}=\dfrac{10}{12}=\dfrac{5}{6}$ και $\dfrac{2}{5}\cdot(-2)=-\dfrac{4}{5}$.\par
\elegxos $\frac{5}{6}\cdot\frac{4}{5}=\frac{20}{30}=\frac{2}{3}$.

\paradeigma{Σύνθετα κλάσματα}
\tixero Τα $\dfrac{\frac{3}{4}}{\frac{9}{8}}$, $\dfrac{\frac{5}{6}}{10}$ και $\dfrac{4}{\frac{2}{3}}$.\par
\tithelo Να τα γράψω ως απλά κλάσματα.\par
\possindeo Γράφω κάθε ένα ως διαίρεση: $\frac{3}{4}:\frac{9}{8}=\frac{3}{4}\cdot\frac{8}{9}=\frac{24}{36}=\frac{2}{3}$. Το δεύτερο: $\frac{5}{6}:10=\frac{5}{6}\cdot\frac{1}{10}=\frac{5}{60}=\frac{1}{12}$. Το τρίτο: $4:\frac{2}{3}=4\cdot\frac{3}{2}=6$.\par
\apantisi $\dfrac{2}{3}$, $\dfrac{1}{12}$, $6$.\par
\elegxos Στο τρίτο: $6\cdot\frac{2}{3}=4$.

\paradeigma{Παράσταση με σύνθετο κλάσμα}
\tixero Η παράσταση $\dfrac{\frac{3}{4}+\frac{1}{2}}{\frac{5}{6}-\frac{1}{3}}$.\par
\tithelo Να βρω την τιμή της.\par
\possindeo Εκτελώ τις πράξεις στον αριθμητή και στον παρονομαστή: $\frac{3}{4}+\frac{1}{2}=\frac{5}{4}$ και $\frac{5}{6}-\frac{1}{3}=\frac{5}{6}-\frac{2}{6}=\frac{3}{6}=\frac{1}{2}$. Τότε $\frac{5}{4}:\frac{1}{2}=\frac{5}{4}\cdot2=\frac{10}{4}$.\par
\apantisi $\dfrac{5}{2}$.\par
\elegxos Με δεκαδικούς: $1{,}25:0{,}5=2{,}5$.

\paradeigma{Διαίρεση δεκαδικών}
\tixero Η διαίρεση $3{,}6:0{,}45$ (κιλά αλεύρι προς κιλά ανά κέικ).\par
\tithelo Πόσα κέικ βγαίνουν;\par
\possindeo Πολλαπλασιάζω και τους δύο αριθμούς επί $100$, ώστε ο διαιρέτης να γίνει φυσικός: $360:45=8$.\par
\apantisi $8$ κέικ.\par
\elegxos $8\cdot0{,}45=3{,}6$.

\paradeigma{Διαίρεση με $10$, $100$, $0{,}1$, $0{,}01$}
\tixero Ο αριθμός $12{,}68$.\par
\tithelo Να τον διαιρέσω με $10$, $100$, $0{,}1$, $0{,}01$.\par
\possindeo Με $10$ και $100$ η υποδιαστολή πάει αριστερά κατά $1$ και $2$ θέσεις. Με $0{,}1$ και $0{,}01$ πάει δεξιά κατά $1$ και $2$ θέσεις.\par
\apantisi $1{,}268$, \ $0{,}1268$, \ $126{,}8$, \ $1\,268$.\par
\elegxos Το $12{,}68:0{,}1$ είναι «πόσα δέκατα χωράνε στο $12{,}68$», δηλαδή $126{,}8$.

\tmima{Ασκήσεις}

\zoni{Το κατάλαβα}{κανόνες και πράξεις}

\ex Κάνε τις διαιρέσεις: α) $15:(-5)$ \quad β) $(-60):(+15)$ \quad γ) $(+81):(-9)$ \quad δ) $(-490):(-7)$.

\ex Βρες τα πηλίκα: α) $\frac{2}{5}:\frac{4}{15}$ \quad β) $\left(+\frac{1}{2}\right):\left(-\frac{1}{3}\right)$ \quad γ) $4:\frac{1}{2}$ \quad δ) $\left(-\frac{4}{3}\right):(-8)$ \quad ε) $\frac{18}{5}:\frac{6}{5}$ \quad στ) $\frac{12}{21}:\frac{7}{4}$.

\ex Μετάτρεψε τα σύνθετα κλάσματα σε απλά: α) $\dfrac{\frac{7}{11}}{\frac{3}{22}}$ \quad β) $\dfrac{\frac{27}{49}}{\frac{9}{14}}$ \quad γ) $\dfrac{4}{\frac{56}{9}}$ \quad δ) $\dfrac{\frac{11}{2}}{121}$ \quad ε) $\dfrac{\frac{3}{4}}{2\frac{1}{8}}$.

\ex Υπολόγισε τα πηλίκα: α) $4{,}5:0{,}15$ \quad β) $28{,}42:49$ \quad γ) $35{,}76:1{,}5$ \quad δ) $292{,}62:29{,}262$ \quad ε) $0{,}48:0{,}6$.

\ex Υπολόγισε χωρίς κάθετη πράξη: α) $45{,}7:10$ \quad β) $3{,}2:100$ \quad γ) $0{,}5:1\,000$ \quad δ) $6{,}4:0{,}1$ \quad ε) $0{,}35:0{,}01$ \quad στ) $12:0{,}001$.

\ex Σωστό ή λάθος; α) $\alpha:\beta=\alpha\cdot\frac{1}{\beta}$. \quad β) Επιτρέπεται η διαίρεση με το $0$. \quad γ) Το πηλίκο ετεροσήμων είναι αρνητικό. \quad δ) $(-12):(-4)=-3$. \quad ε) $3:\frac{1}{3}=1$. \quad στ) Το πηλίκο δύο ρητών είναι πάντα μικρότερο από τον διαιρετέο.

\ex Υπολόγισε: α) $\dfrac{-3}{4}-\dfrac{5}{-8}+\dfrac{-1}{-2}$ \quad β) $1+\dfrac{5}{-4}-\dfrac{-1}{8}$.

\ex Βρες τον αριθμό $x$ ώστε: α) $(-3)\cdot x=12$ \quad β) $x\cdot\frac{2}{3}=4$ \quad γ) $(-5)\cdot x=-\frac{1}{2}$.

\zoni{Το χρησιμοποιώ}{μοιράζω και μετράω}

\ex \textbf{Το λάδι.} Ο παππούς Γιάννης έχει $6$ λίτρα λάδι και θέλει να το βάλει σε μπουκάλια των $\frac{3}{4}$ του λίτρου. Πόσα μπουκάλια θα γεμίσουν;

\ex \textbf{Η κορδέλα.} Η Αγγελική κόβει μια κορδέλα $4{,}5$ μ. σε κομμάτια των $0{,}15$ μ. Πόσα κομμάτια βγαίνουν;

\ex \textbf{Η τιμή.} Η γιαγιά Μυρτώ πλήρωσε $6$ € για $0{,}75$ κιλά τυρί. Πόσο κοστίζει το κιλό;

\ex \textbf{Το χρέος του ταμείου.} Το ταμείο της εκδρομής έχει υπόλοιπο $-24{,}60$ €. Τα $6$ παιδιά της παρέας μοιράζονται εξίσου το χρέος. Τι αριθμός εκφράζει το μερίδιο του καθενός;

\ex \textbf{Η ταχύτητα.} Ο Νίκος τρέχει $\frac{3}{5}$ του χιλιομέτρου σε $\frac{1}{10}$ της ώρας. Πόσα χιλιόμετρα τρέχει σε μία ώρα (αν διατηρεί τον ίδιο ρυθμό);

\zoni{Το σκέφτομαι}{γιατί το πηλίκο μεγαλώνει;}

\ex Ο Βασίλης λέει ότι η διαίρεση πάντα μικραίνει τον αριθμό. Βρες αντιπαράδειγμα και εξήγησε πότε το πηλίκο είναι μεγαλύτερο από τον διαιρετέο.

\ex Γιατί δεν επιτρέπεται η διαίρεση με το $0$; Σκέψου τι θα σήμαινε $5:0=x$ αν υπήρχε τέτοιος αριθμός.

\ex Ποιο είναι μεγαλύτερο: το $3:\frac{1}{2}$ ή το $3:\frac{1}{3}$; Εξήγησε με τη φράση «πόσες φορές χωράει».

\ex Αν $\frac{\alpha}{\beta}>0$, τι πρόσημα μπορούν να έχουν οι $\alpha$ και $\beta$; Αν $\frac{\alpha}{\beta}<0$;

\zoni{Για το σπίτι}{κόβω και μετράω}

\ex Πάρε ένα σπάγκο μήκους $1{,}2$ μ. και κόψε τον σε κομμάτια των $0{,}15$ μ. Πρόβλεψε πρώτα με διαίρεση πόσα θα βγουν και μετά έλεγξε.

\ex Γράψε τρία σύνθετα κλάσματα της επιλογής σου και μετάτρεψέ τα σε απλά.

\paragrafos{Δυνάμεις ρητών: οι πολλαπλασιασμοί που επαναλαμβάνονται}{Α.3.6, «Δυνάμεις ρητών με εκθέτη μη μηδενικό φυσικό»}

\begin{anagke}
Στην αυλή του σχολείου ο Νίκος αφήνει μια μπάλα του μπάσκετ από ύψος $2$ μέτρων. Κάθε φορά που αναπηδά, φτάνει στο μισό του ύψους της προηγούμενης αναπήδησης. «Σε τι ύψος θα φτάσει στην τρίτη αναπήδηση;» ρωτάει η Πολυξένη. Ο Ηλίας γράφει: μετά την πρώτη $2\cdot\frac{1}{2}$, μετά τη δεύτερη $2\cdot\frac{1}{2}\cdot\frac{1}{2}$, μετά την τρίτη $2\cdot\frac{1}{2}\cdot\frac{1}{2}\cdot\frac{1}{2}$. «Αυτό γράφεται πιο σύντομα», λέει η Ρήνη: «το $\frac{1}{2}$ τρεις φορές». Στους φυσικούς το ξέρουν ήδη: $2\cdot2\cdot2=2^3$. Μπορεί να γίνει και με κλάσματα;

Η κυρία Αλεξάνδρα βάζει στον πίνακα δύο παραστάσεις: $\left(-\frac{1}{2}\right)^3$ και $-\frac{1}{2^3}$. «Είναι ίδιες;» ρωτά. Ο Στέλιος λέει «ναι, όλα είναι μείον ένα όγδοο» και η Φανή γράφει στο τετράδιό της: «ίδια, αλλά γιατί;». Ο Βασίλης προτείνει $(-3)^2=-9$ και η Ζωή τον ρωτάει «πώς βγαίνει αρνητικό το γινόμενο δύο αρνητικών;»

Σε αυτή την παράγραφο θα επεκτείνουμε τις δυνάμεις στους ρητούς: θα μάθουμε να προσδιορίζουμε το πρόσημο μιας δύναμης και να υπολογίζουμε παραστάσεις με δυνάμεις.
\end{anagke}

\begin{lexi}{δύναμη, βάση, εκθέτης}{δύναμη, βάση, εκθέτης}
Η \lexiword{δύναμη} είναι η «ύψωση», το γινόμενο ίσων παραγόντων. Στα αρχαία ελληνικά «δύναμις» ήταν και το τετράγωνο. Η \lexiword{βάση} είναι ο αριθμός που πολλαπλασιάζεται με τον εαυτό του και ο \lexiword{εκθέτης} (από το «εκ-τίθημι», «βάζω έξω») ο αριθμός που γράφουμε πάνω δεξιά και δείχνει πόσους παράγοντες έχουμε.
\end{lexi}

\begin{metafrasi}
Το $\frac{2}{3}$ τρεις φορές. & $\left(\dfrac{2}{3}\right)^{3}=\dfrac{2}{3}\cdot\dfrac{2}{3}\cdot\dfrac{2}{3}=\dfrac{8}{27}$\\
Δύο αρνητικοί παράγοντες. & $(-1{,}5)^2=+2{,}25$\\
Τρεις αρνητικοί παράγοντες. & $(-0{,}1)^3=-0{,}001$\\
Η δύναμη με εκθέτη $1$. & $\alpha^1=\alpha$\\
Το μείον μόνο, χωρίς παρένθεση. & $-2^4=-16$\\
\end{metafrasi}

\tmima{Η θεωρία}

\begin{orismos}{δύναμη ρητού}
Η \textbf{δύναμη} του $\alpha$ στην $\nu$ ($\nu$ θετικός φυσικός) είναι το γινόμενο $\nu$ παραγόντων ίσων με $\alpha$:
$$\alpha^{\nu}=\underbrace{\alpha\cdot\alpha\cdot\ldots\cdot\alpha}_{\nu\ \text{παράγοντες}}.$$
Ο $\alpha$ λέγεται \textbf{βάση} (είναι ρητός) και ο $\nu$ \textbf{εκθέτης}. Ισχύει $\alpha^1=\alpha$.
\end{orismos}

Για παράδειγμα $\left(\frac{2}{3}\right)^3=\frac{2}{3}\cdot\frac{2}{3}\cdot\frac{2}{3}=\frac{2^3}{3^3}=\frac{8}{27}$. Γενικά, αφού το γινόμενο κλασμάτων έχει αριθμητή το γινόμενο των αριθμητών και παρονομαστή το γινόμενο των παρονομαστών, ισχύει $\left(\dfrac{\alpha}{\beta}\right)^{\nu}=\dfrac{\alpha^{\nu}}{\beta^{\nu}}$.

\begin{kanonas}{το πρόσημο μιας δύναμης}
\begin{itemize}[leftmargin=1.2em,itemsep=0.3ex,topsep=2pt]
\item Δύναμη με βάση \textbf{θετικό} αριθμό είναι πάντα \textbf{θετική}: αν $\alpha>0$, τότε $\alpha^{\nu}>0$.
\item Δύναμη με βάση \textbf{αρνητικό} αριθμό και εκθέτη \textbf{άρτιο} είναι \textbf{θετική}: αν $\alpha<0$ και $\nu$ άρτιος, τότε $\alpha^{\nu}>0$. Π.χ. $(-1{,}5)^2=(-1{,}5)\cdot(-1{,}5)=+2{,}25$.
\item Δύναμη με βάση \textbf{αρνητικό} αριθμό και εκθέτη \textbf{περιττό} είναι \textbf{αρνητική}: αν $\alpha<0$ και $\nu$ περιττός, τότε $\alpha^{\nu}<0$. Π.χ. $(-0{,}1)^3=-0{,}001$.
\end{itemize}
\end{kanonas}

Ο λόγος είναι ο κανόνας των προσήμων του γινομένου: κάθε ζευγάρι αρνητικών παραγόντων δίνει θετικό, άρα το πρόσημο εξαρτάται από το αν ο αριθμός των αρνητικών παραγόντων είναι άρτιος ή περιττός. Ειδικά $(-1)^{\nu}=1$ αν ο $\nu$ είναι άρτιος και $-1$ αν είναι περιττός. Επίσης $1^{\nu}=1$ και $0^{\nu}=0$.

\tmimasmall{Προτεραιότητα των πράξεων}

Στις παραστάσεις με πράξεις εκτελούμε με τη σειρά: πρώτα τις πράξεις μέσα σε \textbf{παρενθέσεις}, μετά τις \textbf{δυνάμεις}, μετά τους \textbf{πολλαπλασιασμούς και τις διαιρέσεις} (από αριστερά προς τα δεξιά) και τέλος τις \textbf{προσθέσεις και αφαιρέσεις}.

\begin{prosoxi}
\begin{itemize}[leftmargin=1.2em,itemsep=0.3ex,topsep=2pt]
\item Το $\alpha^{3}$ δεν είναι $3\cdot\alpha$, αλλά $\alpha\cdot\alpha\cdot\alpha$. Για παράδειγμα $0{,}2^4\neq4\cdot0{,}2$: $0{,}2^4=0{,}0016$.
\item Το $(-3)^2$ και το $-3^2$ δεν είναι ίδια. Το πρώτο είναι $(-3)(-3)=9$ και το δεύτερο είναι ο αντίθετος του $3^2$, δηλαδή $-9$. Η βάση είναι ό,τι βρίσκεται μέσα στην παρένθεση.
\item Στο $\left(\frac{2}{3}\right)^3$ ο εκθέτης δουλεύει και στον αριθμητή και στον παρονομαστή: $\frac{2^3}{3^3}=\frac{8}{27}$ και όχι $\frac{2^3}{3}$.
\item Ο εκθέτης δείχνει πλήθος παραγόντων, όχι πρόσθεση: $2^3=2\cdot2\cdot2=8$, όχι $2+2+2=6$.
\item Η δύναμη με άρτιο εκθέτη δεν είναι ποτέ αρνητική, ακόμη και με αρνητική βάση.
\end{itemize}
\end{prosoxi}

\tmima{Λυμένα παραδείγματα}

\paradeigma{Η μπάλα του Νίκου}
\tixero Η μπάλα αφήνεται από ύψος $2$ μ. και κάθε αναπήδηση φτάνει στο μισό της προηγούμενης.\par
\tithelo Σε τι ύψος φτάνει στην τρίτη αναπήδηση;\par
\possindeo Κάθε αναπήδηση πολλαπλασιάζει το ύψος επί $\frac{1}{2}$. Μετά από τρεις: $2\cdot\left(\frac{1}{2}\right)^3=2\cdot\frac{1}{8}=\frac{2}{8}=\frac{1}{4}$.\par
\apantisi $\frac{1}{4}$ μ., δηλαδή $25$ εκατοστά.\par
\elegxos Βήμα βήμα: $2\to1\to\frac{1}{2}\to\frac{1}{4}$.

\paradeigma{Δυνάμεις με πρόσημο}
\tixero Οι δυνάμεις $\left(\frac{2}{3}\right)^3$, $(-1{,}5)^2$, $(-0{,}1)^3$.\par
\tithelo Να τις υπολογίσω.\par
\possindeo Πρώτα το πρόσημο: $\frac{2}{3}$ θετική βάση, άρα θετικό. Το $(-1{,}5)^2$ έχει αρνητική βάση και άρτιο εκθέτη, άρα θετικό. Το $(-0{,}1)^3$ έχει αρνητική βάση και περιττό εκθέτη, άρα αρνητικό. Μετά οι απόλυτες τιμές: $\frac{8}{27}$, $1{,}5\cdot1{,}5=2{,}25$, $0{,}1\cdot0{,}1\cdot0{,}1=0{,}001$.\par
\apantisi $\dfrac{8}{27}$, $\ 2{,}25$, $\ -0{,}001$.\par
\elegxos Το $0{,}001$ έχει τρία δεκαδικά ψηφία: $1+1+1$.

\paradeigma{Η παρένθεση μετράει}
\tixero Οι $(-2)^4$, $-2^4$, $\left(-\frac{1}{2}\right)^3$ και $-\frac{1}{2^3}$.\par
\tithelo Να τις υπολογίσω και να τις συγκρίνω.\par
\possindeo Το $(-2)^4$ έχει αρνητική βάση και άρτιο εκθέτη: $(-2)\cdot(-2)\cdot(-2)\cdot(-2)=16$. Το $-2^4$ είναι ο αντίθετος του $2^4$: $-16$. Το $\left(-\frac{1}{2}\right)^3=-\frac{1}{8}$. Το $-\frac{1}{2^3}=-\frac{1}{8}$.\par
\apantisi $16$, $-16$, $-\dfrac{1}{8}$, $-\dfrac{1}{8}$. Οι δύο τελευταίες είναι ίσες.\par
\elegxos Το $(-2)^4$ είναι θετικό, ενώ το $-2^4$ αρνητικό: διαφέρουν.

\paradeigma{Η ειδική περίπτωση του $-1$}
\tixero Οι δυνάμεις $(-1)^{100}$, $(-1)^{101}$, $1^{50}$, $0^3$.\par
\tithelo Να τις υπολογίσω χωρίς να γράψω όλους τους παράγοντες.\par
\possindeo Το $(-1)^{\nu}$ είναι $1$ για άρτιο και $-1$ για περιττό $\nu$. Το $1^{\nu}=1$ και $0^{\nu}=0$.\par
\apantisi $1$, $-1$, $1$, $0$.\par
\elegxos Σε ζεύγη: $(-1)^{100}=\big((-1)(-1)\big)^{50}=1^{50}=1$.

\paradeigma{Παράσταση με προτεραιότητα}
\tixero Η παράσταση $3\cdot\left(\frac{1}{2}-\frac{5}{6}\right)+\left(\frac{1}{3}-\frac{2}{3}\right)^2$.\par
\tithelo Να βρω την τιμή της.\par
\possindeo Πρώτα οι παρενθέσεις: $\frac{1}{2}-\frac{5}{6}=\frac{3}{6}-\frac{5}{6}=-\frac{1}{3}$ και $\frac{1}{3}-\frac{2}{3}=-\frac{1}{3}$. Μετά η δύναμη: $\left(-\frac{1}{3}\right)^2=\frac{1}{9}$. Μετά ο πολλαπλασιασμός: $3\cdot\left(-\frac{1}{3}\right)=-1$. Τέλος η πρόσθεση: $-1+\frac{1}{9}=-\frac{9}{9}+\frac{1}{9}=-\frac{8}{9}$.\par
\apantisi $-\dfrac{8}{9}$.\par
\elegxos Με δεκαδικούς: $3\cdot(-0{,}333\ldots)+0{,}111\ldots=-0{,}888\ldots$

\tmima{Ασκήσεις}

\zoni{Το κατάλαβα}{ορισμός, πρόσημο, υπολογισμός}

\ex Σωστό ή λάθος; α) $\left(\frac{5}{2}\right)^2=\frac{5^2}{2^2}$. \quad β) $\left(\frac{2}{3}\right)^3=\frac{2^3}{3}$. \quad γ) $0{,}2^4=4\cdot0{,}2$. \quad δ) $\left(-\frac{3}{10}\right)^2=0{,}3^2$. \quad ε) $\left(-\frac{1}{3}\right)^3=-\left(\frac{1}{3}\right)^3$.

\ex Υπολόγισε: α) $\left(\frac{1}{2}\right)^2$ \quad β) $\left(-\frac{1}{3}\right)^3$ \quad γ) $\left(\frac{2}{5}\right)^2$ \quad δ) $(-0{,}2)^2$ \quad ε) $0{,}3^3$ \quad στ) $(-1)^7$.

\ex Βρες το πρόσημο χωρίς να υπολογίσεις: α) $\left(-\frac{3}{4}\right)^5$ \quad β) $(-0{,}2)^4$ \quad γ) $\left(-\frac{1}{2}\right)^{10}$ \quad δ) $-\left(\frac{3}{5}\right)^3$ \quad ε) $(-7)^{11}$.

\ex Γράψε ως δύναμη: α) $\frac{3}{5}\cdot\frac{3}{5}\cdot\frac{3}{5}$ \quad β) $(-2)(-2)(-2)(-2)$ \quad γ) $0{,}1\cdot0{,}1$ \quad δ) $\left(-\frac{1}{2}\right)\cdot\left(-\frac{1}{2}\right)$.

\ex Υπολόγισε και σύγκρινε: α) $(-2)^4$ και $-2^4$ \quad β) $(-3)^3$ και $-3^3$ \quad γ) $\left(-\frac{1}{2}\right)^2$ και $-\left(\frac{1}{2}\right)^2$.

\ex Υπολόγισε: α) $1^{100}$ \quad β) $0^5$ \quad γ) $(-1)^{100}$ \quad δ) $(-1)^{101}$ \quad ε) $10^3$ \quad στ) $(0{,}1)^2$.

\ex Κάνε τις πράξεις: α) $\left(\frac{1}{2}\right)^2+\left(-\frac{1}{2}\right)^3+\frac{1}{4^2}$ \quad β) $0{,}1^2+\left(\frac{1}{10}\right)^2-0{,}2^2$ \quad γ) $-1^{10}+1^{100}+(-1)^2+(-1)^3$.

\zoni{Το χρησιμοποιώ}{μπάλες, διπλώσεις, εμβαδά}

\ex \textbf{Η μπάλα.} Ο Νίκος αφήνει μια μπάλα από ύψος $2$ μ. και κάθε αναπήδηση φτάνει στο μισό της προηγούμενης. Σε τι ύψος φτάνει μετά την τρίτη και μετά την τέταρτη αναπήδηση;

\ex \textbf{Οι διπλώσεις.} Η Αγγελική διπλώνει ένα φύλλο χαρτί στη μέση πέντε φορές. Τι μέρος του αρχικού εμβαδού έχει το διπλωμένο φύλλο; Γράψε το ως δύναμη και ως κλάσμα.

\ex \textbf{Το τετράγωνο.} Μια αυλή είναι τετράγωνο με πλευρά $1{,}5$ μ. Πόσο είναι το εμβαδόν της; Ο Ηλίας λέει $1{,}5\cdot2=3$. Τι λάθος έκανε;

\ex \textbf{Η παράσταση.} Υπολόγισε $8\cdot\left(-\frac{1}{2}\right)^3+(0{,}5+0{,}3)^2$.

\ex \textbf{Το μοτίβο.} Υπολόγισε τα $\left(\frac{1}{10}\right)^1$, $\left(\frac{1}{10}\right)^2$, $\left(\frac{1}{10}\right)^3$, $\left(\frac{1}{10}\right)^4$ ως δεκαδικούς. Τι παρατηρείς;

\zoni{Το σκέφτομαι}{γιατί βγαίνει έτσι;}

\ex Ο Βασίλης υπολόγισε $-3^2=9$. Εξήγησε το λάθος του και δείξε ότι $(-3)^2=9$.

\ex Πότε είναι θετική μια δύναμη με αρνητική βάση; Εξήγησε με τον κανόνα των προσήμων.

\ex Ποιο είναι μεγαλύτερο, το $\left(\frac{1}{2}\right)^3$ ή το $\left(\frac{1}{2}\right)^2$; Και στη βάση $2$; Τι παρατηρείς για βάσεις μικρότερες και μεγαλύτερες από το $1$;

\ex Αν $\alpha=-\frac{1}{2}$, ποιο είναι μεγαλύτερο, το $\alpha^2$ ή το $\alpha^3$; Εξήγησε χωρίς πράξεις.

\zoni{Για το σπίτι}{πίνακες δυνάμεων}

\ex Φτιάξε πίνακα με τις δυνάμεις $\left(\frac{1}{2}\right)^{\nu}$ για $\nu=1,2,\dots,6$, με κλάσμα και δεκαδικό. Τι παρατηρείς;

\ex Φτιάξε τρεις παραστάσεις με ρητούς και δυνάμεις, που να χρειάζονται και παρενθέσεις, και υπολόγισέ τες.

\paragrafos{Τυποποιημένη μορφή: οι μεγάλοι αριθμοί χωράνε σε μια γραμμή}{Α.3.7, «Τυποποιημένη μορφή μεγάλων αριθμών»}

\begin{anagke}
Η κυρία Αλεξάνδρα ζητάει από την Α3 μια εργασία για το ηλιακό σύστημα. Η Αγγελική γράφει στο τετράδιό της την απόσταση της Γης από τον Ήλιο: $149\,600\,000\,000$ μέτρα. Ο Ηλίας γράφει την ηλικία της Γης: $4\,500\,000\,000$ έτη, και την ξαναγράφει τρεις φορές, γιατί δεν είναι σίγουρος για τα μηδενικά. Ο Στέλιος δοκιμάζει την ηλικία του Γαλαξία και χάνεται στο δεύτερο μέτρημα. «Αυτό δεν διαβάζεται», παραπονιέται η Ρήνη.

Η Πολυξένη, που έχει το κομπιουτεράκι της, πατάει $149\,600\,000\cdot1\,000$ και η οθόνη δείχνει $1{,}496\mathrm{E}11$. «Τι είναι αυτό το E11;». Η κυρία Αλεξάνδρα χαμογελάει: «Το κομπιουτεράκι δεν χωράει όλα τα μηδενικά και τα γράφει με συντομία: $1{,}496$ επί δέκα στην ενδέκατη. Κάνουν ό,τι κάνουν και οι επιστήμονες». Η Φανή ρωτάει: «Μπορούμε να γράφουμε κι εμείς έτσι, ώστε να μη χάνουμε τα μηδενικά;».

Σε αυτή την παράγραφο θα μάθουμε να γράφουμε τους πολύ μεγάλους αριθμούς με σύντομο και ενιαίο τρόπο, τον οποίο χρησιμοποιούν οι επιστήμονες.
\end{anagke}

\begin{lexi}{τυποποιημένη μορφή}{τυποποιημένη μορφή}
Ο \lexiword{τύπος} σημαίνει καλούπι, πρότυπο. \lexiword{Τυποποιημένη} είναι η μορφή που έχει το ίδιο καλούπι για όλους τους αριθμούς: ένας αριθμός από το $1$ μέχρι κάτω από το $10$, επί μια δύναμη του $10$. Οι επιστήμονες την προτιμούν γιατί φαίνεται αμέσως το \emph{μέγεθος} του αριθμού (πόσα μηδενικά) και τα σημαντικά του ψηφία.
\end{lexi}

\begin{metafrasi}
Ένα ακολουθούμενο από $7$ μηδενικά. & $10^7=10\,000\,000$\\
Τριάντα εκατομμύρια. & $30\,000\,000=3\cdot10^7$\\
Η απόσταση Γης -- Ήλιου σε χιλιόμετρα. & $149\,600\,000=1{,}496\cdot10^8$\\
Στην οθόνη του κομπιουτεράκι. & $1{,}496\mathrm{E}8$\\
\end{metafrasi}

\tmima{Η θεωρία}

Θυμόμαστε ότι η δύναμη $10^{\nu}$ είναι ο αριθμός που γράφεται με το $1$ και μετά $\nu$ μηδενικά: $10^3=1\,000$, $10^6=1\,000\,000$.

\begin{orismos}{τυποποιημένη μορφή}
Ένας «μεγάλος» αριθμός μπορεί να γραφεί στη μορφή $\alpha\cdot10^{\nu}$, δηλαδή ως γινόμενο ενός αριθμού $\alpha$ επί μια δύναμη του $10$, όπου ο $\nu$ είναι ένας μη μηδενικός φυσικός αριθμός. Η μορφή αυτή λέγεται \textbf{τυποποιημένη}. Ο $\alpha$ είναι ένας δεκαδικός αριθμός με ακέραιο μέρος ένα ψηφίο μεγαλύτερο του $0$, δηλαδή $1\le\alpha<10$.
\end{orismos}

Η τυποποιημένη μορφή μας δίνει τη δυνατότητα να γράψουμε έναν «μεγάλο» αριθμό με απλούστερο τρόπο.

\begin{kanonas}{από τον αριθμό στην τυποποιημένη μορφή}
\begin{enumerate}[leftmargin=1.6em,itemsep=0.2ex,topsep=2pt]
\item Βάζω την υποδιαστολή μετά το \emph{πρώτο ψηφίο} (αριστερά), ώστε να προκύψει αριθμός $\alpha$ από $1$ έως $9{,}99\ldots$
\item Μετράω πόσες θέσεις μετακίνησα την υποδιαστολή προς τα αριστερά. Αυτός είναι ο εκθέτης $\nu$.
\item Γράφω $\alpha\cdot10^{\nu}$ (και αφαιρώ τα μηδενικά που περισσεύουν στο τέλος του $\alpha$).
\end{enumerate}
Για παράδειγμα $815\,000\,000\,000=8{,}15\cdot10^{11}$, γιατί η υποδιαστολή μετακινήθηκε $11$ θέσεις προς τα αριστερά. Ο εκθέτης είναι πάντα κατά $1$ μικρότερος από το πλήθος των ψηφίων του αριθμού.
\end{kanonas}

\begin{kanonas}{από την τυποποιημένη μορφή στον αριθμό}
Για να γράψω τον $\alpha\cdot10^{\nu}$ ως συνηθισμένο αριθμό, μεταφέρω την υποδιαστολή του $\alpha$ \textbf{$\nu$ θέσεις προς τα δεξιά}, συμπληρώνοντας με μηδενικά όπου χρειάζεται. Έτσι $3{,}42\cdot10^9=3\,420\,000\,000$.
\end{kanonas}

Η τυποποιημένη μορφή βοηθά και στη σύγκριση. Από δύο αριθμούς σε τυποποιημένη μορφή μεγαλύτερος είναι αυτός με τον \textbf{μεγαλύτερο εκθέτη}. Αν έχουν τον ίδιο εκθέτη, συγκρίνουμε τους $\alpha$. Έτσι $2\cdot10^{5}>9{,}9\cdot10^{4}$ και $3{,}2\cdot10^{6}<4{,}1\cdot10^{6}$.

\begin{prosoxi}
\begin{itemize}[leftmargin=1.2em,itemsep=0.3ex,topsep=2pt]
\item Το $12\cdot10^{5}$ δεν είναι τυποποιημένη μορφή, γιατί το $12$ δεν είναι μικρότερο από το $10$. Σωστά: $1{,}2\cdot10^{6}$.
\item Το $0{,}5\cdot10^{3}$ δεν είναι τυποποιημένη, γιατί το ακέραιο μέρος είναι $0$. Σωστά: $5\cdot10^{2}$.
\item Το $10^{\nu}$ δεν είναι $10\cdot\nu$. Το $10^{3}=1\,000$ και όχι $30$.
\item Ο εκθέτης δεν είναι το πλήθος των ψηφίων αλλά το πλήθος των μηδενικών μετά το $1$, δηλαδή τα ψηφία μείον $1$: ο $5\,000\,000$ έχει $7$ ψηφία και είναι $5\cdot10^{6}$.
\item Τα μηδενικά στο τέλος του $\alpha$ δεν γράφονται: $2{,}50\cdot10^4=2{,}5\cdot10^4$.
\end{itemize}
\end{prosoxi}

\tmima{Λυμένα παραδείγματα}

\paradeigma{Στην τυποποιημένη μορφή}
\tixero Οι αριθμοί $30\,000\,000$ και $815\,000\,000\,000$.\par
\tithelo Να τους γράψω σε τυποποιημένη μορφή.\par
\possindeo Ο πρώτος έχει $8$ ψηφία: μετακινώ την υποδιαστολή κατά $7$ θέσεις και παίρνω $3$. Ο δεύτερος έχει $12$ ψηφία: κατά $11$ θέσεις, και παίρνω $8{,}15$.\par
\apantisi $30\,000\,000=3\cdot10^{7}$ και $815\,000\,000\,000=8{,}15\cdot10^{11}$.\par
\elegxos $3\cdot10^{7}=3\cdot10\,000\,000$ και $8{,}15\cdot10^{11}$: η υποδιαστολή πάει $11$ θέσεις δεξιά, και παίρνουμε το $815$ και $9$ μηδενικά.

\paradeigma{Από την τυποποιημένη μορφή}
\tixero Ο αριθμός $3{,}42\cdot10^{9}$.\par
\tithelo Να τον γράψω ως συνηθισμένο αριθμό.\par
\possindeo Μετακινώ την υποδιαστολή κατά $9$ θέσεις δεξιά. Το $3{,}42$ έχει δύο δεκαδικά ψηφία, οπότε προσθέτω $7$ μηδενικά.\par
\apantisi $3\,420\,000\,000$.\par
\elegxos Ο αριθμός έχει $10$ ψηφία, ενώ ο εκθέτης είναι $9$.

\paradeigma{Σύγκριση}
\tixero Οι αριθμοί $2\cdot10^{3}$ και $3\cdot10^{2}$, και οι $0{,}1\cdot10^{2}$ και $10$.\par
\tithelo Να τους συγκρίνω.\par
\possindeo Στο πρώτο ζεύγος οι εκθέτες διαφέρουν: $2\,000>300$. Το δεύτερο ζεύγος: το $0{,}1\cdot10^{2}$ δεν είναι τυποποιημένο. Είναι $0{,}1\cdot100=10$.\par
\apantisi $2\cdot10^{3}>3\cdot10^{2}$ και $0{,}1\cdot10^{2}=10$.\par
\elegxos $2\cdot10^3=2\,000$ και $3\cdot10^2=300$.

\paradeigma{Χρόνος σε λεπτά}
\tixero Οι καρχαρίες της Γροιλανδίας ζουν περίπου $400$ χρόνια. Ένα έτος έχει περίπου $365$ ημέρες.\par
\tithelo Πόσα λεπτά είναι τα $400$ χρόνια; Να δώσω την απάντηση σε τυποποιημένη μορφή.\par
\possindeo Ένα έτος έχει $365\cdot24\cdot60=525\,600$ λεπτά. Τα $400$ έτη έχουν $400\cdot525\,600=210\,240\,000$ λεπτά. Ο αριθμός έχει $9$ ψηφία, άρα ο εκθέτης είναι $8$.\par
\apantisi $2{,}1024\cdot10^{8}$ λεπτά.\par
\elegxos $2{,}1024\cdot10^8$: η υποδιαστολή πάει $8$ θέσεις δεξιά και δίνει $210\,240\,000$.

\paradeigma{Ποιοι αριθμοί είναι τυποποιημένοι;}
\tixero Οι $4{,}5\cdot10^{3}$, $45\cdot10^{2}$, $0{,}45\cdot10^{4}$, $1\cdot10^{6}$.\par
\tithelo Να βρω ποιοι είναι σε τυποποιημένη μορφή και να διορθώσω τους άλλους.\par
\possindeo Ελέγχω αν το ακέραιο μέρος του πρώτου παράγοντα είναι ένα ψηφίο από το $1$ ως το $9$. Στο $4{,}5$ και στο $1$ ναι. Το $45$ έχει δύο ψηφία και το $0{,}45$ έχει ακέραιο μέρος $0$.\par
\apantisi Τυποποιημένοι είναι οι $4{,}5\cdot10^{3}$ και $1\cdot10^{6}$. Οι άλλοι δύο είναι ίσοι με τον πρώτο, $45\cdot10^2=4\,500=0{,}45\cdot10^4=4{,}5\cdot10^3$.\par
\elegxos Και οι τρεις δίνουν $4\,500$.

\tmima{Ασκήσεις}

\zoni{Το κατάλαβα}{από και προς την τυποποιημένη μορφή}

\ex Γράψε σε τυποποιημένη μορφή: α) $150\,000$ \quad β) $21\,000$ \quad γ) $1\,300\,000$ \quad δ) $2\,030\,000$.

\ex Γράψε σε πλήρη μορφή: α) $2{,}31\cdot10^{4}$ \quad β) $1{,}36\cdot10^{9}$ \quad γ) $1{,}37\cdot10^{2}$ \quad δ) $10^{6}$.

\ex Σύγκρινε τους αριθμούς με ένα από τα $<,\ >,\ =$: α) $2\cdot10^{3}\ \ldots\ 3\cdot10^{2}$ \quad β) $0{,}1\cdot10^{2}\ \ldots\ 10$ \quad γ) $1{,}7\cdot10^{4}\ \ldots\ 3\cdot10^{3}$ \quad δ) $6\cdot10^{5}\ \ldots\ 0{,}6\cdot10^{6}$.

\ex Ποιοι από τους αριθμούς $4{,}5\cdot10^{3}$, $45\cdot10^{2}$, $0{,}45\cdot10^{4}$, $9{,}99\cdot10^{8}$, $10\cdot10^{5}$, $1\cdot10^{6}$ είναι σε τυποποιημένη μορφή;

\ex Γράψε σε τυποποιημένη μορφή τους αριθμούς που δεν είναι, από την προηγούμενη άσκηση.

\ex Γράψε στη μορφή $\alpha\cdot10^{\nu}$ τους αριθμούς: α) $7\,000$ \quad β) $480\,000\,000$ \quad γ) $12\,500\,000\,000$ \quad δ) $6\,021\,000$.

\ex Σωστό ή λάθος; α) $3\cdot10^{5}=300\,000$. \quad β) $2{,}5\cdot10^{4}=250\,000$. \quad γ) Ο $10^{7}$ γράφεται με το $1$ και επτά μηδενικά. \quad δ) $5{,}2\cdot10^{3}>4{,}9\cdot10^{3}$. \quad ε) $3\cdot10^{2}=6$. \quad στ) $7\cdot10^{1}=70$.

\zoni{Το χρησιμοποιώ}{διάστημα, χρόνος, πληθυσμοί}

\ex \textbf{Τα αστρονομικά μεγέθη.} Γράψε σε τυποποιημένη μορφή: α) Απόσταση Γης -- Σελήνης: $384\,400$ χιλιόμετρα. \quad β) Απόσταση Γης -- Ήλιου: $149\,600\,000$ χιλιόμετρα. \quad γ) Ηλικία της Γης: $4\,500\,000\,000$ έτη. \quad δ) Ηλικία του Γαλαξία: $13\,600\,000\,000$ έτη. \quad ε) Απόσταση Πλούτωνα -- Ήλιου: $5\,900\,000\,000$ χιλιόμετρα.

\ex \textbf{Η χελώνα.} Μια γιγάντια χελώνα ζει περίπου $150$ χρόνια. Σε πόσα λεπτά αντιστοιχούν (με έτος $365$ ημερών); Γράψε την απάντηση σε τυποποιημένη μορφή.

\ex \textbf{Οι πλανήτες.} Οι αποστάσεις από τον Ήλιο: Δίας $7{,}8\cdot10^{8}$ km, Άρης $2{,}3\cdot10^{8}$ km, Ποσειδώνας $4{,}5\cdot10^{9}$ km, Ερμής $5{,}8\cdot10^{7}$ km. Γράψε τους πλανήτες από τον πιο κοντινό στον πιο μακρινό.

\ex \textbf{Το κομπιουτεράκι.} Το κομπιουτεράκι της Πολυξένης δείχνει $3{,}2\mathrm{E}8$ και μετά $4{,}05\mathrm{E}6$. Ποιοι είναι οι αριθμοί;

\ex \textbf{Οι πληθυσμοί.} Η Κίνα έχει περίπου $1\,400\,000\,000$ κατοίκους και η Ελλάδα περίπου $10\,400\,000$. Γράψε τους αριθμούς σε τυποποιημένη μορφή και πες ποιος είναι μεγαλύτερος και πώς το βλέπεις από τους εκθέτες.

\zoni{Το σκέφτομαι}{γιατί; και τι θα συνέβαινε αν...;}

\ex Εξήγησε γιατί το $12\cdot10^{5}$ δεν είναι σε τυποποιημένη μορφή και γράψε το σωστά.

\ex Πόσα ψηφία έχει ένας αριθμός που γράφεται στη μορφή $\alpha\cdot10^{\nu}$; Πόσα ψηφία έχουν οι $4\cdot10^{9}$ και $4{,}5\cdot10^{9}$;

\ex Ο Στέλιος έγραψε $5\,400\,000=5{,}4\cdot10^{5}$. Βρες το λάθος του.

\ex Ανάμεσα σε ποιες δύο διαδοχικές δυνάμεις του $10$ βρίσκεται ο $4{,}5\cdot10^{6}$;

\zoni{Για το σπίτι}{ψάξε και γράψε}

\ex Ψάξε πέντε μεγάλα μεγέθη από την εγκυκλοπαίδεια ή το διαδίκτυο (π.χ. πληθυσμός της Γης, μάζα ή ακτίνα πλανήτη) και γράψε τα σε τυποποιημένη μορφή.

\ex Υπολόγισε πόσα δευτερόλεπτα έχει ένα έτος ($365$ ημέρες) και γράψε το αποτέλεσμα σε τυποποιημένη μορφή.

\epanalipsi{Α3}

\tmima{Λεξικό του κεφαλαίου}

\renewcommand{\arraystretch}{1.5}
\begin{longtable}{@{}>{\raggedright\arraybackslash}p{0.28\linewidth} >{\raggedright\arraybackslash}p{0.62\linewidth}@{}}
\toprule
\textbf{Λέξη ή τύπος} & \textbf{Με μία πρόταση}\\
\midrule
\endhead
ρητοί & Οι αριθμοί που γράφονται ως κλάσμα $\frac{\alpha}{\beta}$ με $\alpha,\beta$ ακέραιους και $\beta\neq0$.\\
ισοδύναμα κλάσματα & Ίδια αξία: $\frac{\alpha}{\beta}=\frac{\gamma}{\delta}$ όταν $\alpha\delta=\beta\gamma$.\\
ανάγωγο & Κλάσμα που δεν απλοποιείται: $\MKD$ των όρων ίσος με $1$.\\
$|\alpha|$ & Η απόσταση του $\alpha$ από το $0$: ποτέ αρνητική.\\
αντίθετοι & Ετερόσημοι με ίδια απόλυτη τιμή: $\alpha$ και $-\alpha$, με άθροισμα $0$.\\
σύγκριση & Μεγαλύτερος είναι ο δεξιότερος στον άξονα.\\
$\alpha-\beta$ & $=\alpha+(-\beta)$: προσθέτω τον αντίθετο.\\
γινόμενο, πηλίκο & Ομόσημοι: θετικό. Ετερόσημοι: αρνητικό. Ποτέ διαίρεση με το $0$.\\
αντίστροφος & Ο $\frac{\beta}{\alpha}$ για τον $\frac{\alpha}{\beta}$: γινόμενο $1$. Το $0$ δεν έχει.\\
$\alpha^{\nu}$ & $\nu$ παράγοντες ίσοι με $\alpha$. Ίδιο πρόσημο με το γινόμενο.\\
$\alpha\cdot10^{\nu}$ & Τυποποιημένη μορφή: $1\le\alpha<10$ και $\nu$ μη μηδενικός φυσικός.\\
\bottomrule
\end{longtable}

\tmima{Ασκήσεις επανάληψης}

\ex Για καθέναν από τους $-\frac{6}{3}$, $\frac{4}{5}$, $0$, $-2{,}5$, $\frac{7}{7}$ πες αν είναι φυσικός, ακέραιος, ρητός.

\ex Απλοποίησε το $\frac{36}{48}$ μέχρι να γίνει ανάγωγο και γράψε ένα κλάσμα ισοδύναμο με το $\frac{5}{6}$ που έχει παρονομαστή $42$.

\ex Γράψε το $\frac{7}{8}$ ως δεκαδικό και τον $0{,}36$ ως ανάγωγο κλάσμα.

\ex Βρες: α) $|-7{,}5|$ \quad β) $\left|\frac{3}{4}\right|$ \quad γ) τον αντίθετο του $-\frac{2}{3}$ \quad δ) όλους τους $\alpha$ με $|\alpha|=4$.

\ex Γράψε σε αύξουσα σειρά τους αριθμούς $-2{,}5$, $0$, $-\frac{1}{2}$, $\frac{3}{2}$, $-3$, $1{,}2$.

\ex Βρες δύο ρητούς ανάμεσα στους $\frac{2}{7}$ και $\frac{3}{7}$.

\ex Υπολόγισε: α) $-\frac{3}{4}+\frac{5}{6}$ \quad β) $2-\frac{7}{5}$ \quad γ) $-1{,}5+0{,}75$ \quad δ) $\frac{1}{2}-\left(-\frac{1}{3}\right)$.

\ex Υπολόγισε την παράσταση $-\frac{1}{2}+\frac{2}{3}-\frac{1}{6}+\frac{1}{4}$.

\ex Υπολόγισε: α) $\left(-\frac{3}{4}\right)\cdot\frac{8}{9}$ \quad β) $(-0{,}5)\cdot(-1{,}2)$ \quad γ) $12\cdot\left(\frac{1}{3}+\frac{1}{4}\right)$.

\ex Υπολόγισε: α) $\left(-\frac{3}{5}\right):\frac{9}{10}$ \quad β) $6:\left(-\frac{3}{4}\right)$ \quad γ) $7{,}2:0{,}09$.

\ex Μετάτρεψε σε απλό κλάσμα το $\dfrac{\frac{2}{3}+\frac{1}{6}}{\frac{5}{4}}$.

\ex Υπολόγισε: α) $\left(-\frac{2}{3}\right)^3$ \quad β) $-3^2$ \quad γ) $(-0{,}1)^2$ \quad δ) $\frac{1}{2}+\left(-\frac{1}{2}\right)^2$.

\ex Γράψε σε τυποποιημένη μορφή τον $45\,000\,000$, σε πλήρη μορφή τον $3{,}02\cdot10^{5}$ και σύγκρινε τους $8\cdot10^{6}$ και $9{,}5\cdot10^{5}$.

\ex \textbf{Ο αγώνας.} Ένα σχολείο έχει $240$ μαθητές. Τα $\frac{3}{8}$ παίζουν μπάσκετ και το $\frac{1}{3}$ αυτών είναι κορίτσια. Πόσα κορίτσια παίζουν μπάσκετ;

\ex \textbf{Σωστό ή λάθος (πράξεις).} α) $-\frac{3}{4}<-\frac{2}{3}$. \quad β) $\frac{1}{3}+\frac{1}{3}=\frac{2}{6}$. \quad γ) $(-2)\cdot(-3)\cdot(-4)=-24$. \quad δ) $5:\frac{1}{5}=1$. \quad ε) Το $0$ έχει αντίστροφο. \quad στ) Ο αντίθετος του $0$ είναι το $0$.

\ex \textbf{Σωστό ή λάθος (δυνάμεις, τυποποιημένη μορφή).} α) $(-2)^3=-6$. \quad β) $-2^2=-4$. \quad γ) $\left(\frac{2}{3}\right)^2=\frac{4}{3}$. \quad δ) $5\cdot10^{4}=50\,000$. \quad ε) Το $12\cdot10^{3}$ είναι σε τυποποιημένη μορφή. \quad στ) Τα $\frac{3}{4}$ και $\frac{9}{12}$ είναι ισοδύναμα.

\Needspace{24\baselineskip}\tmima{Ο χάρτης της επανάληψης}

Μετά τις λύσεις, δες σε ποιες ασκήσεις δυσκολεύτηκες. Ο πίνακας σου λέει ποια παράγραφο να ξαναδιαβάσεις.

\renewcommand{\arraystretch}{1.5}
\begin{tabular*}{\linewidth}{@{\extracolsep{\fill}} >{\raggedright\arraybackslash}p{0.38\linewidth} >{\raggedright\arraybackslash}p{0.52\linewidth}}
\toprule
\textbf{Αν δυσκολεύτηκες στις} & \textbf{ξαναδιάβασε την παράγραφο}\\
\midrule
ΕΑ3.1, ΕΑ3.2, ΕΑ3.3 & Α3.1 (ρητοί, ισοδύναμα, δεκαδικοί)\\
ΕΑ3.4, ΕΑ3.5, ΕΑ3.6 & Α3.2 (απόλυτη τιμή, διάταξη, ενδιάμεσοι ρητοί)\\
ΕΑ3.7, ΕΑ3.8 & Α3.3 (πρόσθεση και αφαίρεση)\\
ΕΑ3.9, ΕΑ3.14 & Α3.4 (πολλαπλασιασμός, κλάσμα ποσότητας)\\
ΕΑ3.10, ΕΑ3.11 & Α3.5 (διαίρεση, σύνθετα κλάσματα)\\
ΕΑ3.12 & Α3.6 (δυνάμεις)\\
ΕΑ3.13 & Α3.7 (τυποποιημένη μορφή)\\
ΕΑ3.15, ΕΑ3.16 & Α3.1 έως Α3.7\\
\bottomrule
\end{tabular*}

